\documentclass[11pt]{article}
\usepackage[utf8]{inputenc}
\usepackage[T1]{fontenc}
\usepackage{lmodern}
\usepackage[margin=1in]{geometry}
\usepackage{amsmath,amssymb,amsthm,mathtools}
\usepackage{booktabs,longtable,array,tabularx}
\usepackage{enumitem}
\usepackage{microtype}
\usepackage{xcolor}
\usepackage{hyperref}
\usepackage{fancyhdr}
\usepackage{titlesec}
\usepackage{tocloft}
\usepackage{mdframed}
\usepackage{url}
\usepackage{ragged2e}

\hypersetup{colorlinks=false,pdfborder={0 0 0},pdftitle={An AASC Constraint-Formalism Proof of the Poincare Endpoint by Fixed-Carrier Negative-Branch Exclusion},pdfauthor={Amos Jay Maley},pdfsubject={AASC constraint-formalism proof of the Poincare endpoint},pdfkeywords={Poincare Conjecture, AASC, endpoint under audit, negative branch exclusion, fixed carrier}}
\pagestyle{fancy}
\setlength{\headheight}{14pt}
\fancyhf{}
\lhead{AASC Poincare Closure}
\rhead{Negative-Branch Bridge Audit}
\cfoot{\thepage}
\setlist[itemize]{topsep=3pt,itemsep=2pt,parsep=0pt,leftmargin=2.0em}
\setlist[enumerate]{topsep=3pt,itemsep=2pt,parsep=0pt,leftmargin=2.2em}
\renewcommand{\arraystretch}{1.16}
\setlength{\LTpre}{6pt}
\setlength{\LTpost}{6pt}
\setlength{\parskip}{0.35em}
\setlength{\parindent}{0pt}
\emergencystretch=3em
\hbadness=10000
\newcolumntype{Y}{>{\RaggedRight\arraybackslash}X}
\newcolumntype{L}[1]{>{\RaggedRight\arraybackslash}p{#1}}

\newtheorem{theorem}{Theorem}[section]
\newtheorem{lemma}[theorem]{Lemma}
\newtheorem{proposition}[theorem]{Proposition}
\newtheorem{corollary}[theorem]{Corollary}
\newtheorem{definition}[theorem]{Definition}
\newtheorem{remark}[theorem]{Remark}
\newtheorem{audit}[theorem]{Audit Node}
\newtheorem{nonclaim}[theorem]{Nonclaim}

\newcommand{\AASC}{\mathrm{AASC}}
\newcommand{\ATS}{\mathrm{ATS}}
\newcommand{\UEAP}{\mathrm{UEAP}}
\newcommand{\Sthree}{S^{3}}
\newcommand{\ClosedThree}{\mathrm{Closed3}}
\newcommand{\Conn}{\mathrm{Conn}}
\newcommand{\StandPC}{\mathrm{Stand}_{PC}}
\newcommand{\SphereRead}{\mathrm{SphereRead}}
\newcommand{\NonSphere}{\mathrm{NonSphere}}
\newcommand{\CarrierPC}{C_{PC}}
\newcommand{\EndpointAdeqPC}{\mathrm{EndpointAdeq}_{PC}}
\newcommand{\KPC}{\mathcal{K}_{PC}}
\newcommand{\PCEndpoint}{\mathrm{PCEndpointUnderAudit}}
\newcommand{\OfficialPCEndpointUse}{\mathrm{OfficialPCEndpointUse}}
\newcommand{\PCCarrierInst}{\mathrm{PCCarrierInstantiated}}
\newcommand{\OfficialPCNegativeResolution}{\mathrm{OfficialPCNegativeResolution}}
\newcommand{\PCNeg}{\mathrm{PCNeg}}
\newcommand{\StdPCNeg}{\mathrm{StdPCNeg}}
\newcommand{\SphereBridgeImgExcl}{\mathrm{SphereBridgeImgExcl}_{PC}}
\newcommand{\ThmSphereDisc}{\mathrm{ThmSphereDisc}_{PC}}
\newcommand{\EndpointGov}{\mathrm{EndpointGov}_{PC}}
\newcommand{\IndependentSphereDisc}{\mathrm{IndependentSphereDisc}_{PC}}
\newcommand{\BridgeSlotPC}{B^{\mathrm{slot}}_{PC}}
\newcommand{\BridgeCompPC}{B^{\mathrm{comp}}_{PC}}
\newcommand{\BridgeObjPC}{B^{PC}}
\newcommand{\BridgeObjCompPC}{B^{PC,\mathrm{comp}}}
\newcommand{\SphereBridgeImgExclLoc}{\mathrm{SphereBridgeImgExcl}^{\mathrm{loc}}_{PC}}
\newcommand{\LocalCountercasePC}{\mathrm{LocalCountercase}_{PC}}
\newcommand{\EndpointCounterforcePC}{\mathrm{EndpointCounterforce}_{PC}}
\newcommand{\ReductioCountercasePC}{\mathrm{ReductioCountercase}_{PC}}
\newcommand{\Homeo}{\cong_{\mathrm{homeo}}}
\newcommand{\Diffeo}{\cong_{\mathrm{diff}}}
\newcommand{\Disc}{\mathrm{Disc}}
\newcommand{\Adm}{\mathrm{Adm}}
\newcommand{\PhiPC}{\Phi_{PC}}
\newcommand{\Meta}{\mathrm{Meta}}

\title{\textbf{An AASC Constraint-Formalism Proof of the Poincare Endpoint by Fixed-Carrier Negative-Branch Exclusion}\\[0.4em]
\large Kernel-First Endpoint Use, Sphere-Bridge Image Exclusion, No-Fifth-Case Governance, and Route-Locus Exhaustion}
\author{Amos Jay Maley\\Independent Researcher, Ang Thong, Thailand}
\date{June 2026}

\begin{document}
\pagenumbering{roman}
\maketitle

\begin{center}
\fbox{\parbox{0.92\linewidth}{\textbf{Proof class.} This manuscript is a mathematical proof in the AASC constraint formalism. It is not a conventional Ricci-flow proof, surgery proof, triangulation proof, recognition algorithm proof, or imported corollary of the Hamilton--Perelman theorem. Its proof class is kernel-forced exclusion of the native Poincare negative branch on the fixed closed simply connected three-manifold carrier. The manuscript theorem chain is the proof. Lean 4 material is support/audit material; the current Poincare Lean audit archive records the endpoint-spine support surface, and the stable paper-specific formalization scope includes the relevant AASC machinery directly.}}
\end{center}

\begin{abstract}
This manuscript gives an audit-grade mathematical proof in the AASC constraint formalism of the Poincare Conjecture by fixed-carrier exclusion of the native negative endpoint branch. The proof is independent of the Hamilton--Perelman Ricci-flow-with-surgery theorem, Morgan--Tian, Kleiner--Lott, Cao--Zhu, or any equivalent classical solution of Poincare as a premise. Those sources enter only after the AASC endpoint, as comparators for proof burden and bridge architecture.

AASC is used here as mathematics: a constraint formalism with theorem-bearing primitives, kernel roles, endpoint-governance rules, and no-independent-classifier closure. The proof is the manuscript theorem chain. Lean 4 support, where referenced or included, audits formal routing and kernel discipline. The current Poincare Lean audit archive records the endpoint-spine support surface, and the stable paper-specific formalization scope includes the relevant AASC machinery directly. Lean is not substituted for this manuscript proof.

The proof does not infer an admissible discriminator directly from true non-spherehood. Instead, it uses the official Poincare theorem as a determinate endpoint under audit and fixes the pointwise carrier
\[
\CarrierPC(M)=(M,\ClosedThree,\Conn,\dim M=3,\pi_1(M),\Sthree,\Homeo).
\]
Endpoint-bearing use of this theorem already requires target determinacy, step evaluability, act-time finality, and same-regime fidelity. These kernel-neutral endpoint-adequacy conditions force the local Poincare kernel roles: reference, standing, admissibility, irreversibility, fail-closed bivalence, standing--admissibility identity, no generators, no carriers, no repairs, no selector authority, UEAP report preservation, ATS skin non-authority, and no second same-domain endpoint classifier. The live K-nodes that exclude independent negative governance are also audited for weakening resistance: any proposed weaker same-carrier regime either preserves the same endpoint behavior extensionally, lowers the branch to support or bookkeeping, shifts carrier, or admits forbidden second endpoint authority.

The negative branch is routed in stages. First, the native counterexample normal form is fixed:
\[
\PCNeg(M)\Longleftrightarrow \ClosedThree(M)\wedge \Conn(M)\wedge \pi_1(M)=0\wedge M\not\Homeo\Sthree.
\]
This branch is then identified with sphere-bridge image exclusion on the fixed Poincare carrier. When used as the official negative resolution of the pointwise endpoint, sphere-bridge image exclusion is theorem-level endpoint-status governance. Endpoint-status governance that does not occupy the positive sphere bridge induces an independent same-domain sphere discriminator. Such a discriminator is excluded by the endpoint-forced kernel: it is either positive bridge completion, proof support only, carrier/domain shift, bookkeeping, or forbidden second endpoint authority. Hence \(\PCNeg(M)\) is excluded for every fixed closed connected simply connected 3-manifold \(M\), and therefore \(M\Homeo\Sthree\).
\end{abstract}

\noindent\textbf{Keywords.} Poincare Conjecture; AASC; ATS; UEAP; fixed carrier; endpoint under audit; negative branch exclusion; sphere-bridge image exclusion; no fifth case; independent discriminator; route-locus exhaustion; Perelman; Ricci flow.

\section*{Mathematical-status lock and Lean support boundary}
\addcontentsline{toc}{section}{Mathematical-status lock and Lean support boundary}

\begin{mdframed}[backgroundcolor=gray!8,linewidth=0.6pt]
\textbf{Mathematical-status lock.} AASC is used here as a mathematical constraint formalism, not as philosophical commentary, metaphor, optional narration, or an interpretive overlay.  The proof is the theorem chain stated in this manuscript.  The AASC primitives used here---reference, standing, admissibility, irreversibility, bivalence, no-generator closure, no-carrier transfer, no-repair, AMetric no-selector discipline, UEAP report preservation, ATS layer discipline, and no-independent-classifier closure---are theorem-bearing mathematical constraints.  A same-mode objection must therefore address the mathematical constraint chain, the kernel-to-Poincare instantiation, or the endpoint-governance inferences.  It is not a refutation to classify AASC as non-mathematical vocabulary or to demand a Ricci-flow-native proof route.
\end{mdframed}

\begin{mdframed}[backgroundcolor=gray!8,linewidth=0.6pt]
\textbf{Lean support boundary.} The proof of this paper is the manuscript proof, not a Lean substitute.  Lean 4 material is kept to a support appendix and proof-routing map in this release.  The public Poincare Lean audit archive records the endpoint-spine support surface for the manuscript-facing route, while the stable paper-specific formalization scope includes the relevant AASC machinery directly: kernel discipline, below-kernel underivability posture, endpoint routing, and no-independent-classifier closure.  A challenge to Lean support must identify a mismatch between the cited support surface and the manuscript theorem chain.  It is not enough to observe that Lean is not itself the proof, because the manuscript does not claim that it is.
\end{mdframed}

\section*{Read this first: endpoint proof spine}
\addcontentsline{toc}{section}{Read this first: endpoint proof spine}

The live proof spine below is the manuscript proof spine in the AASC mathematical constraint formalism. It is not a classical topological construction, and it is not replaced by Lean. Lean-facing material, where mentioned, is audit support for this chain.

It is intentionally short.
\[
\begin{aligned}
&\text{official Poincare endpoint use}\Rightarrow \text{non-degenerate fixed Poincare carrier}\Rightarrow \KPC,\\
&\PCNeg(M)\Longleftrightarrow \StdPCNeg(M)\Longleftrightarrow \SphereBridgeImgExcl(M),\\
&\SphereBridgeImgExcl(M)\wedge \OfficialPCNegativeResolution(M)\Rightarrow \ThmSphereDisc(M),\\
&\ThmSphereDisc(M)\Rightarrow \EndpointGov(M)\Rightarrow \IndependentSphereDisc(M)\Rightarrow \bot.
\end{aligned}
\]
Before the negative branch is used, the local packet is audited for weakening resistance. In particular, K5, K6, K11, and K13 are shown to be local endpoint necessities rather than optional reinforcements: weakening them while preserving determinate same-carrier endpoint use either leaves the endpoint behavior unchanged, lowers the alleged negative branch to support or bookkeeping, shifts carrier, or creates forbidden second endpoint authority.

Therefore \(\neg\PCNeg(M)\). Since \(\PCNeg(M)\) is the ordinary pointwise counterexample normal form, \(\neg\PCNeg(M)\) is exactly the ordinary sphere-readout conclusion under the fixed standing packet:
\[
\ClosedThree(M)\wedge\Conn(M)\wedge\pi_1(M)=0\Longrightarrow M\Homeo\Sthree.
\]
Consequently, a same-mode objection must break one of the mathematical links in the displayed chain. A request for Ricci flow, surgery, or a conventional 3-manifold construction is a request for a different proof class, not a refutation of this proof spine.

\begin{mdframed}[backgroundcolor=gray!8,linewidth=0.6pt]
\textbf{No one-step retyping.} The proof does not say ``true non-spherehood supplies an admissible discriminator.'' The native negative branch first appears as standard counterexample normal form. It becomes endpoint-status governance only if it is used as the official negative resolution of the fixed pointwise endpoint. This is the central safeguard against collapsing the native negative branch into its own discriminator.
\end{mdframed}

\begin{mdframed}[backgroundcolor=gray!8,linewidth=0.6pt]
\textbf{Local exact-complement route.} The global negative-resolution route remains the route audit for official endpoint-defeating use. The endpoint closeout can also be run locally. Under fixed \(\StandPC(M)\), the local assumption \([M\not\Homeo\Sthree]_i\), equivalently \([\neg\SphereRead(M)]_i\), is opened as the live exact-complement countercase against sphere-readout. It is not promoted to global official negative endpoint resolution. It receives only local endpoint-countercase standing inside the subproof, so the induced contradiction discharges the local non-spherehood assumption itself.
\[
\begin{aligned}
[\neg\SphereRead(M)]_i
&\Rightarrow \LocalCountercasePC(M)
 \Rightarrow \EndpointCounterforcePC(M),\\
&\Rightarrow \SphereBridgeImgExclLoc(M)
 \Rightarrow \IndependentSphereDisc(M)
 \Rightarrow \bot.
\end{aligned}
\]
\end{mdframed}

\subsection*{Dependency inversion exclusion}
The dependency order is not
\[
\text{AASC as optional overlay}\Rightarrow\text{Poincare constraints}.
\]
The dependency order is
\[
\begin{aligned}
&\text{Poincare as theorem-bearing endpoint}\\
&\Rightarrow\text{non-degenerate fixed-carrier regime}\\
&\Rightarrow\text{kernel roles already active}.
\end{aligned}
\]
AASC names and audits roles already required for a determinate same-carrier theorem endpoint. It does not impose those roles from outside topology as optional vocabulary. The names are AASC names; the constraints are mathematical necessities of non-degenerate fixed-domain endpoint use.

\subsection*{Cost of denying the kernel}
A denial of \(\KPC\) is not a free refusal of vocabulary. It abandons at least one minimal condition required for the Poincare endpoint to function as a determinate theorem-bearing target.

\begin{longtable}{L{0.23\linewidth}L{0.67\linewidth}}
\toprule
\textbf{Denied role} & \textbf{Cost in Poincare endpoint terms}\\
\midrule
Reference & The endpoint no longer fixes the same manifold \(M\), the same homeomorphism relation, the same standing packet, or the same sphere-readout slot.\\
Standing & Closedness, connectedness, dimension three, simple connectedness, bridge status, or negative endpoint use no longer carry theorem-bearing status.\\
Admissibility & There is no governed distinction between lawful same-carrier endpoint continuation and carrier drift, surrogate readout, repair, selector import, or second-classifier governance.\\
Irreversibility & Later surgery, flow-stage data, triangulation, decomposition, or category reinterpretation may retroactively repair an earlier unbridged endpoint report as the same act.\\
Bivalence / fail-closed status & Unknown, deferred, unsupported, silent, partial, or unclassified status becomes a third endpoint truth branch on the same fixed carrier.\\
Standing--admissibility identity & A branch may carry endpoint standing while inadmissible, or admissibility may perform endpoint work without standing.\\
No second classifier & A negative-status discriminator may govern the same sphere endpoint beside the declared bridge while preserving the appearance of one fixed carrier.\\
\bottomrule
\end{longtable}

\subsection*{What must be denied to reject the proof}
A same-mode rejection must deny at least one live theorem link.
\begin{longtable}{L{0.36\linewidth}L{0.56\linewidth}}
\toprule
\textbf{Proof link} & \textbf{Standing denial burden}\\
\midrule
Endpoint under audit is target fixation & Show that naming the official universal Poincare target already assumes \(M\Homeo\Sthree\) for a pointwise carrier.\\
Pointwise endpoint use & Show that an arbitrary closed simply connected 3-manifold under the universal theorem is not evaluated under the pointwise endpoint slot.\\
Endpoint use instantiates the kernel & Exhibit a theorem-bearing same-carrier Poincare endpoint act that preserves reference, standing, admissibility, and irreversibility while lacking the kernel roles.\\
Native negative normal form & Produce a Poincare counterexample not expressible as closed, connected, simply connected, three-dimensional, and not homeomorphic to \(\Sthree\).\\
Bridge correspondence & Show that the native negative normal form is not sphere-bridge image exclusion on the fixed carrier.\\
Endpoint-bearing bridge exclusion & Show official negative use of fixed-carrier sphere-bridge exclusion that does not classify endpoint status.\\
No hidden fifth case & Inhabit endpoint-resolving, fixed-carrier, non-governing negative endpoint use.\\
Endpoint governance implies independent discriminator & Show endpoint-resolving negative governance that is positive bridge completion, proof support only, carrier shift, bookkeeping, or otherwise not second endpoint authority.\\
No independent sphere discriminator & Show an independent same-domain sphere-status discriminator compatible with the endpoint-forced kernel.\\
Strict weakening of the local kernel packet & Exhibit a weaker same-carrier regime that preserves target-slot fixation, minimal report evaluability, standing--admissibility discipline, and carrier fidelity while permitting endpoint-resolving negative governance that is not support, bookkeeping, carrier shift, bridge completion, coequal target role, or second endpoint authority.\\
Sphere-bridge object & Show that a clause of \(\BridgeObjPC(M)\) smuggles \(\SphereRead(M)\), changes carrier, or fails to preserve the official Poincare endpoint interface.\\
Local exact-complement discharge & Show that the local annotation is an added object-level premise or that the contradiction discharges only a strengthened premise rather than \([\neg\SphereRead(M)]_i\).\\
Official endpoint correspondence & Show that excluding \(\PCNeg(M)\) under the fixed standing packet does not yield \(M\Homeo\Sthree\).\\
AASC mathematical status & Show a precise failure of a definition, theorem, inference, dependency claim, or formal constraint; not merely that AASC is unfamiliar or nonclassical.\\
Lean support boundary & Show a mismatch between Lean-facing support and a manuscript claim, if Lean support is invoked; not merely that Lean is not the proof.\\
\bottomrule
\end{longtable}

\tableofcontents
\newpage
\pagenumbering{arabic}

\section{Proof-Class, Claim-Class, and Source-Use Locks}

\subsection{Proof-class lock}
This manuscript is not a conventional 3-manifold topology proof in the Hamilton--Perelman style. It does not run Ricci flow, perform surgery, prove canonical-neighborhood estimates, prove finite extinction, or import geometrization. Those routes may be legitimate bridge-construction proofs in the classical language. The present proof is a mathematical proof in the AASC constraint formalism: it excludes the native pointwise counterexample branch once that branch is used as official negative endpoint resolution. Its proof obligations are kernel validity, fixed-carrier endpoint adequacy, negative-branch use classification, endpoint-governance closure, no-independent-discriminator exclusion, and official endpoint correspondence.

\begin{mdframed}[backgroundcolor=gray!8,linewidth=0.6pt]
\textbf{Proof-Class Lock.} The proof spine is
\[
\begin{aligned}
\PCNeg&\to\SphereBridgeImgExcl\to\OfficialPCNegativeResolution\\
&\to\EndpointGov\to\IndependentSphereDisc\to\bot.
\end{aligned}
\]
The live proof does not depend on resemblance to Ricci flow, surgery, or customary 3-manifold proof presentation.
\end{mdframed}

\subsection{Claim-class lock}
The manuscript claims:
\begin{quote}
AASC proves the Poincare endpoint by showing that official negative use of the native counterexample branch induces endpoint-status governance, that such governance induces an independent same-domain sphere discriminator, and that the fixed-carrier kernel excludes independent same-domain endpoint discrimination. Therefore the native negative branch is impossible and the ordinary sphere-readout endpoint follows.
\end{quote}
It does not claim a new Ricci-flow proof, a new surgery theorem, a new recognition algorithm, an imported corollary of Perelman, or a proof that a topological truth must first present itself as an independently named invariant. It also does not claim that negative statements are forbidden. A negative claim may be proof support, bookkeeping, local evidence, a lawful obstruction inside a bridge route, or a carrier-changing report. What is excluded is a same-carrier theorem-bearing negative endpoint resolution that governs the fixed endpoint without occupying the positive sphere bridge. Lean 4 support is not offered as a replacement proof of this manuscript; it remains an audit-support boundary for the already-mathematical AASC machinery and formal routing, now recorded in the public Poincare Lean audit archive.

\begin{nonclaim}[Classical-route boundary]
The manuscript is not offered as a conventional topology derivation. It is an AASC endpoint-closure proof. A conventional objection has force only if it identifies a failed theorem link in the declared proof spine, not merely because the route is not Ricci flow.
\end{nonclaim}

\subsection{Source-use lock}
Classical sources may be used for problem statement, historical context, category hygiene, and post-proof comparison. They are not proof premises for the live AASC endpoint spine. Moise's theorem may be used for three-dimensional category discipline, not as a Poincare-completion theorem \cite{Moise1977}. Hamilton, Perelman, Morgan--Tian, Kleiner--Lott, and Clay materials are comparison-only sources for the traditional bridge route \cite{Hamilton1982,Perelman2002,Perelman2003Surgery,Perelman2003Extinction,MorganTian2007,KleinerLott2008,ClayPoincare,ClayMorganTian}.

AASC corpus sources supply the mathematical proof-governance kernel that is locally restated below: non-degenerate construction, admissibility, standing--admissibility identity, extensional uniqueness of the admissible interior, bivalence, no generators, no carriers, no repairs, AMetric no-selector boundary, ATS layer discipline, and UEAP report preservation \cite{AASCMatrix2026,PreferredMatrixScan2026,AASC_Kernel_2026,AASC_Structure_Admissibility_2026,AASC_StandingIdentity_2026,AASC_ClaimStanding_UEAP_2026,AASC_AMetric_2026,AASC_ATS_2026,AASC_Bivalence_2026}. The AASC machinery is mathematical in nature; references to Lean support are references to audit support for formal routing and kernel discipline, not to a substitute proof text.

\subsection{Preferred matrix import discipline}
The preferred/hardened AASC matrix rows remain the import pool for kernel support. Superseded rows are not first-line support. The project package preserves the corpus scan, preferred import shortlist, and prior kernel provenance ledger. In this manuscript, those imports are used after endpoint-neutral adequacy has already forced the local role packet. Thus the kernel is not presented as an arbitrary formalism choice; it is a named audit of non-degenerate endpoint use.

\begin{longtable}{L{0.25\linewidth}L{0.32\linewidth}L{0.34\linewidth}}
\toprule
\textbf{Imported result} & \textbf{Preferred source family} & \textbf{Poincare use}\\
\midrule
Extensional uniqueness of admissible interior & Kernel of Admissibility; Structure of Admissibility & Blocks plural independent endpoint-status interiors for the same sphere slot.\\
Standing equals reuse-stable admissibility & Kernel; Standing-Admissibility Identity Closure & Endpoint-resolving negative use cannot carry standing outside admissibility.\\
No second same-domain classifier & Standing-Admissibility Identity Closure; ATS & Excludes independent sphere-status governance beside the bridge.\\
Fixed-domain endpoint exhaustion & Bivalence; unique admissible interior; no deferred admissibility & Blocks unknown, silent, deferred, unsupported, and hidden-fifth status as endpoint branches.\\
No generators/no repair/no carriers & Bivalence; Kernel; Impossibility Suite & Prevents endpoint labels, later repairs, and surrogate carriers from authorizing endpoint status.\\
UEAP report preservation & Claim Standing and Legitimacy & Negative endpoint reports may not exceed fixed carrier, standing, and bridge support.\\
AMetric no-selector & AMetric Boundary & No preferred triangulation, metric, presentation, decomposition, or ranking decides endpoint status.\\
ATS skin non-authority & Anchor, Tensor, and Skin & Expressions, labels, and bookkeeping do not become proof-bearing by being written.\\
Route-locus exhaustion & Same-domain discriminator and carrier-transfer exhaustion rows & Endpoint-resolving negative routes must enter a live route locus or do no endpoint work.\\
\bottomrule
\end{longtable}

\section{Official Poincare Endpoint and Fixed Carrier}\label{sec:official-endpoint}

\subsection{Rank-free endpoint statement}
The official Poincare endpoint is the universal statement
\[
\forall M\,\bigl[\ClosedThree(M)\wedge\Conn(M)\wedge\pi_1(M)=0\Rightarrow M\Homeo\Sthree\bigr],
\]
where \(M\) ranges over topological 3-manifolds unless a smooth or PL presentation is explicitly invoked.

\begin{definition}[Poincare standing]
\[
\StandPC(M)\Longleftrightarrow \ClosedThree(M)\wedge\Conn(M)\wedge \dim M=3\wedge \pi_1(M)=0.
\]
This is standing, not readout.
\end{definition}

\begin{definition}[Sphere-readout]
\[
\SphereRead(M)\Longleftrightarrow M\Homeo\Sthree.
\]
This is the declared endpoint readout.
\end{definition}

\begin{definition}[Poincare endpoint under audit]
\(\PCEndpoint\) means that the official universal Poincare theorem target is fixed for evaluation. It is a target-fixation statement, not an assumption that \(\SphereRead(M)\) holds for any pointwise carrier.
\end{definition}

\begin{definition}[Fixed Poincare carrier]
For a closed connected 3-manifold \(M\), define
\[
\CarrierPC(M)=(M,\ClosedThree,\Conn,\dim M=3,\pi_1(M),\Sthree,\Homeo).
\]
The lawful redescription relation \(\Homeo\) preserves the same topological carrier, standing packet, sphere-readout slot, and negative endpoint slot.
\end{definition}

\begin{definition}[Poincare carrier instantiation]
\(\PCCarrierInst(M)\) means that the fixed pointwise carrier \(\CarrierPC(M)\) is the carrier under endpoint audit: the same manifold \(M\), the same homeomorphism relation, the same closed connected dimension-three carrier, the same simple-connectedness standing slot, and the same sphere-readout slot are in force.
\end{definition}

\begin{definition}[Official Poincare endpoint use]
\(\OfficialPCEndpointUse(M)\) means theorem-bearing use of the Poincare endpoint for the fixed pointwise carrier \(\CarrierPC(M)\). It is not a heuristic, label, partial report, or assertion that the conclusion is already true. It is the act of using the official theorem target as the pointwise endpoint under evaluation.
\end{definition}

\begin{theorem}[Universal endpoint binds pointwise endpoint use]\label{thm:universal-binds-pointwise}
If \(\PCEndpoint\) is the theorem target under audit and \(\StandPC(M)\) holds, then the pointwise instance is evaluated under \(\OfficialPCEndpointUse(M)\) and \(\PCCarrierInst(M)\).
\end{theorem}
\begin{proof}
The universal theorem target is a statement over all closed connected simply connected 3-manifolds. Fixing an arbitrary \(M\) in the antecedent fixes the corresponding pointwise endpoint act: the carrier is \(M\), the standing packet is \(\StandPC(M)\), and the readout slot is \(\SphereRead(M)\). This does not assert \(\SphereRead(M)\). It states only that the pointwise endpoint is under evaluation on the fixed carrier.
\end{proof}

\begin{proposition}[Endpoint use fixes carrier instantiation]\label{prop:endpoint-use-carrier}
\[
\OfficialPCEndpointUse(M)\Rightarrow \PCCarrierInst(M).
\]
\end{proposition}
\begin{proof}
Official endpoint use is theorem-bearing use of the Poincare target on \(\CarrierPC(M)\). Therefore the carrier roles listed in \(\CarrierPC(M)\) are fixed before any bridge, negative branch, route-locus, or comparison report is evaluated.
\end{proof}

\section{Kernel-Neutral Endpoint Adequacy and Dependency Order}\label{sec:kernel-neutral}

This section states the endpoint-adequacy conditions without AASC vocabulary and then shows that they force the local kernel roles. This fixes the dependency order: AASC names the mathematical constraints forced by non-degenerate endpoint use, rather than being externally imposed on topology.

\begin{definition}[Raw Poincare endpoint expression]
A raw Poincare endpoint expression is an inscription, proof trace, obstruction claim, invariant report, geometric comparison, surgery report, triangulation claim, decomposition claim, category transfer, selector output, classifier value, or route statement involving some of the symbols
\[
M,\ \ClosedThree(M),\ \Conn(M),\ \pi_1(M)=0,\ M\Homeo\Sthree,\ M\not\Homeo\Sthree.
\]
A raw expression may be bookkeeping, proof support, heuristic evidence, carrier shift, bridge construction, or endpoint-bearing report. It is not automatically theorem-bearing endpoint content.
\end{definition}

\begin{definition}[Kernel-neutral Poincare endpoint adequacy]
A raw Poincare endpoint expression is endpoint-adequate for the official Poincare target only if the following four conditions are present before the expression is used as theorem-bearing endpoint content.
\begin{enumerate}[label=(A\arabic*)]
\item \textbf{Target determinacy.} The same closed connected 3-manifold \(M\), the same lawful redescription relation \(\Homeo\), the same standing packet \(\StandPC(M)\), and the same sphere-readout role are fixed.
\item \textbf{Step evaluability.} A comparison, invariant, obstruction, bridge, surgery, decomposition, selector, or negative branch is evaluable as advancing, failing, supporting, or leaving that same endpoint rather than merely appearing as notation.
\item \textbf{Act-time finality.} Later triangulation, surgery, flow-stage information, geometric model, decomposition, category translation, or reinterpretation creates a new certified act; it does not retroactively repair an earlier unbridged or failed endpoint report as the same act.
\item \textbf{Same-regime fidelity.} Lawful redescription preserves \(M\), closedness, connectedness, dimension three, simple-connectedness standing, the sphere-readout slot, and endpoint classification. Failure of such preservation is carrier drift or scope reset, not same-endpoint continuation.
\end{enumerate}
\end{definition}

\begin{definition}[Non-degenerate Poincare endpoint regime]
A Poincare endpoint regime is non-degenerate when it has endpoint adequacy, at least one standing-bearing endpoint act, lawful redescription identity, fail-closed treatment of out-of-domain moves, and no post-hoc repair of failed endpoint status as the same act.
\end{definition}

\begin{theorem}[Endpoint adequacy forces the kernel roles]\label{thm:adequacy-forces-kernel}
Every endpoint-adequate Poincare theorem act already requires the governance roles later named reference, standing, admissibility, and irreversibility.
\end{theorem}
\begin{proof}
Target determinacy is reference: without it, the manifold, equivalence relation, standing packet, or endpoint role may drift. Step evaluability is standing: without it, a comparison or obstruction is an inscription rather than theorem-bearing endpoint content. Act-time finality is irreversibility: without it, a failed or unbridged report may be repaired retroactively as the same act. Same-regime fidelity is admissibility: it supplies the boundary between lawful same-endpoint continuation and carrier shift, surrogate substitution, selector import, repair, or second-classifier governance. Therefore the kernel roles are forced by endpoint adequacy before AASC names them.
\end{proof}

\begin{theorem}[Poincare endpoint use instantiates the kernel]\label{thm:endpoint-use-kernel}
For every pointwise carrier \(M\),
\[
\OfficialPCEndpointUse(M)\Rightarrow \KPC(M).
\]
\end{theorem}
\begin{proof}
Official endpoint use is theorem-bearing use of the Poincare target on the fixed carrier. Such use must satisfy target determinacy, step evaluability, act-time finality, and same-regime fidelity. By Theorem \ref{thm:adequacy-forces-kernel}, those conditions instantiate the local roles named by \(\KPC(M)\). The kernel is therefore endpoint-forced, not externally imposed.
\end{proof}

\begin{corollary}[No constraint-formalism opt-out]\label{cor:no-constraint-optout}
A critic cannot preserve Poincare as a determinate same-carrier theorem endpoint while rejecting the kernel roles. To opt out of the kernel, the critic must weaken at least one endpoint-adequacy condition. This is not optional AASC vocabulary; the neutral endpoint-adequacy conditions force the roles that AASC names.
\end{corollary}

\begin{theorem}[Kernel-denial cost for Poincare endpoint use]\label{thm:kernel-denial-cost}
Let \(M\) be under official Poincare endpoint audit. Any denial of \(\KPC(M)\) while preserving official endpoint use must abandon at least one of target determinacy, standing, admissible same-carrier discipline, same-regime fidelity, or act-time finality. Such a denial is not a same-mode counterexample unless it supplies an alternative non-degenerate endpoint regime preserving all of those roles.
\end{theorem}
\begin{proof}
By Theorem \ref{thm:endpoint-use-kernel}, official endpoint use instantiates the kernel because the endpoint act already satisfies the adequacy conditions. Denying the kernel while retaining endpoint use therefore rejects at least one adequacy condition. If none is rejected, the kernel remains in force. If one is rejected, the proposal has left the non-degenerate same-carrier Poincare endpoint regime unless a replacement regime preserving the same roles is supplied.
\end{proof}

\begin{audit}[AASC is not optional vocabulary]\label{audit:not-optional-vocab}
The kernel is not imposed as an optional terminology layer. Target determinacy, step evaluability, act-time finality, and same-regime fidelity force the roles that AASC names reference, standing, irreversibility, and admissibility. The names are AASC names; the constraints are mathematical necessities of non-degenerate fixed-domain endpoint use. Lean support, where invoked, audits this formal routing rather than replacing the manuscript proof.
\end{audit}

\section{Local Poincare Kernel Packet}\label{sec:kernel-packet}

The following packet is the local Poincare form of the endpoint-forced fixed-domain kernel. The full upstream derivations remain in the preferred AASC corpus; this manuscript states the consequences used locally.

\begin{longtable}{L{0.16\linewidth}L{0.32\linewidth}L{0.40\linewidth}}
\toprule
\textbf{Node} & \textbf{General necessity} & \textbf{Local Poincare form}\\
\midrule
K1 & Fixed-domain kernel & A theorem-bearing Poincare endpoint requires fixed carrier, standing, endpoint role, admissibility boundary, and act-time irreversibility.\\
K2 & Reference fixation & A report about \(M\) cannot be licensed by silently replacing \(M\), the equivalence relation, the category, the compactness status, or the endpoint role.\\
K3 & Standing-readout separation & \(\StandPC(M)\) is standing; \(\SphereRead(M)\) is readout. If they were identical, Poincare would be vocabulary rather than endpoint theorem.\\
K4 & Bivalence and fail-closed status & On a fixed endpoint role, a theorem-bearing branch is admitted or not admitted. Unknown, deferred, suggestive, or partial status is not a third endpoint truth value.\\
K5 & Unique admissible interior & The fixed endpoint cannot carry plural independent admissibility interiors for the same sphere slot. A second endpoint-status classifier must collapse, reset carrier, or fail.\\
K6 & Standing--admissibility identity & Endpoint standing cannot float outside admissibility. A negative endpoint report with standing must be admissibly typed; otherwise it is support, bookkeeping, or carrier drift.\\
K7 & No generators & The endpoint classifier cannot generate standing for its own negative value. Writing \(\PhiPC(M)=\NonSphere\) does not self-authorize endpoint truth.\\
K8 & No carriers & Homology data, geometric data, Ricci-flow stages, triangulations, decompositions, or recognition procedures cannot carry sphere-readout to \(M\) without a lawful bridge.\\
K9 & No post-selection repair & Later surgery, flow control, decomposition, or reinterpretation creates a new certified report; it does not repair an earlier unbridged report as the same act.\\
K10 & AMetric no-selector boundary & A preferred triangulation, metric, handlebody, presentation, normal form, or ranking cannot decide endpoint status before admissibility is fixed.\\
K11 & UEAP report preservation & A report may not exceed its target, carrier, standing, and bridge support. Negative endpoint use must have tensor-level endpoint role.\\
K12 & ATS layer discipline & Skin, labels, expressive phrasing, unknown status, silence, and bookkeeping do not become proof-bearing by being written.\\
K13 & Endpoint exhaustion & Once carrier, standing packet, endpoint role, and admissibility boundary are fixed, endpoint-bearing branches are exhausted by positive bridge occupation or governed negative endpoint use; unsupported, deferred, bookkeeping, carrier-shift, and hidden-fifth statuses do not occupy endpoint truth.\\
\bottomrule
\end{longtable}

\begin{theorem}[Local kernel packet K1--K13]\label{thm:local-kernel-packet}
For every official Poincare endpoint use of \(M\), the local kernel consequences K1--K13 are in force on \(\CarrierPC(M)\).
\end{theorem}
\begin{proof}
By Theorem \ref{thm:endpoint-use-kernel}, official endpoint use instantiates reference, standing, admissibility, and irreversibility. K1 records their joint fixed-domain role. K2--K3 are reference and role-separation consequences. K4--K6 are fail-closed, unique-interior, and standing--admissibility consequences. K7--K9 block generation, carrier transfer, and repair. K10 blocks selector authority. K11 records report preservation. K12 blocks skin-to-tensor promotion. K13 is the endpoint exhaustion consequence: after the endpoint coordinates are fixed, a purported branch is bridge occupation, proof support, carrier reset, repair attempt, selector/surrogate import, bookkeeping, or forbidden second-classifier governance. No additional same-domain endpoint occupation remains.
\end{proof}

\begin{remark}[Kernel import discipline]
The kernel packet does not prove Ricci flow estimates, surgery control, triangulation recognition, or a new conventional invariant theorem. It determines what can count as theorem-bearing endpoint use once the Poincare carrier and endpoint slot are fixed. Topological data remain legitimate proof support or bridge completion material; they are blocked only when silently promoted into another endpoint role.
\end{remark}


\section{Weakening-Resistance of the Local Kernel Packet}\label{sec:weakening-resistance}

The remaining same-mode pressure on the proof is not whether reference, standing, admissibility, irreversibility, carrier fidelity, target-slot fixation, and minimal report evaluability are needed at all.  Those are already forced by endpoint adequacy.  The live question is whether a strictly weaker local packet could preserve the same non-degenerate Poincare endpoint while permitting the negative endpoint governance that the proof excludes.  This section answers that question at the local packet nodes that do the most work: K5, K6, K11, and K13.

The claim is local.  The manuscript does not require an assertion that every possible mathematical formalism must contain every AASC node in the same vocabulary.  It requires the following narrower theorem: on the fixed Poincare carrier, under theorem-bearing same-endpoint use, weakening these nodes either becomes extensionally equivalent for endpoint purposes or exits non-degenerate same-carrier endpoint use.

\begin{definition}[Core endpoint adequacy packet]
The core endpoint adequacy packet for the fixed Poincare carrier is
\[
C_0(M)=\{\mathrm{Ref}_M,\mathrm{Stand}_M,\mathrm{Adm}_M,\mathrm{Irr}_M\},
\]
together with target determinacy, step evaluability, act-time finality, same-regime fidelity, and the fixed carrier
\[
\CarrierPC(M)=(M,\ClosedThree,\Conn,\dim M=3,\pi_1(M),\Sthree,\Homeo).
\]
This is the part of the endpoint regime already forced before the full local packet is invoked.
\end{definition}

\begin{definition}[Negative-governance-permissive weakening]
A negative-governance-permissive weakening is a regime \(W\) such that:
\begin{enumerate}[label=(\roman*)]
\item \(W\) preserves the same fixed carrier \(\CarrierPC(M)\);
\item \(W\) preserves theorem-bearing target-slot fixation for the same sphere-readout role;
\item \(W\) preserves minimal report evaluability, same-regime fidelity, and act-time finality: reports remain attached to the fixed target, remain classifiable as theorem-bearing or non-theorem-bearing, and do not silently change carrier;
\item \(W\) preserves the core packet \(C_0(M)\);
\item \(W\) permits an endpoint-resolving negative branch that is not proof support only, not bookkeeping, not carrier shift, not bridge completion, and not a lawful coequal target role.
\end{enumerate}
Endpoint determinacy in this definition means only that the carrier, target slot, and report role are fixed enough for theorem-bearing evaluation.  It does not include the full K13 exhaustion of unknown, deferred, unsupported, bookkeeping, and hidden-fifth statuses.  K13 is proved below as the local theorem that those statuses cannot occupy endpoint truth without degenerating the fixed endpoint regime.  Likewise, minimal report evaluability is weaker than K11: it requires only that reports be attached to the fixed target and classifiable by role.  K11 is the local no-overreporting theorem that endpoint force may not exceed carrier, bridge, and tensor support.

The same-mode challenge is whether such a \(W\) can exist without degenerating the fixed endpoint regime.
\end{definition}

\begin{lemma}[Plural route data are permitted; plural endpoint interiors are not]\label{lem:plural-routes-not-interiors}
K5 does not prohibit plural descriptions, proof routes, invariants, metrics, triangulations, decompositions, Ricci-flow traces, category presentations, or bridge-support structures.  It prohibits only plural endpoint-governing admissible interiors for the same fixed endpoint slot when those interiors assign different endpoint status while preserving the same carrier, standing packet, theorem target, and endpoint role.
\end{lemma}
\begin{proof}
Plural support data may coexist without assigning endpoint status.  They remain proof support, auxiliary carriers, representatives, or bridge material until used as theorem-bearing endpoint authority.  If two alleged interiors assign the same endpoint status, the second is extensionally bookkeeping.  If they assign different endpoint status on the same carrier and same endpoint slot, they supply two governance authorities for one fixed theorem target.  That is not harmless plurality of mathematical data; it is plural endpoint governance.
\end{proof}

\begin{lemma}[Local necessity of K5]\label{lem:k5-local-necessity}
Under \(C_0(M)\), a weakening that permits two distinct endpoint-governing admissible interiors for the same sphere-readout slot either collapses to bookkeeping, splits the domain, or creates a second same-domain endpoint classifier.
\end{lemma}
\begin{proof}
Let \(I_1\) and \(I_2\) be alleged endpoint-governing admissible interiors for the same fixed endpoint slot.  If they agree on all endpoint-bearing acts, they are extensionally identical for endpoint purposes and the second is bookkeeping.  If they disagree while preserving the same carrier, standing packet, endpoint role, and admissibility regime, then one endpoint slot receives two admissible governance verdicts.  The disagreement must be resolved by a higher selector, a domain split, or a second classifier.  A higher selector is endpoint governance above the gate; a domain split exits same-carrier use; a second classifier is exactly the forbidden plural endpoint authority.  Therefore K5 is locally necessary for non-degenerate same-carrier endpoint determinacy.
\end{proof}

\begin{lemma}[Local necessity of K6]\label{lem:k6-local-necessity}
Under \(C_0(M)\), endpoint standing outside admissibility is either support-only, carrier shift, or a second endpoint-status gate.
\end{lemma}
\begin{proof}
Suppose a branch \(B\) has endpoint standing while not being admissibility-bound.  If \(B\) does not affect endpoint status, it is support, bookkeeping, or lower-status report material.  If it affects endpoint status while preserving the same carrier, it performs endpoint-status work outside the admissibility boundary and therefore acts as a second gate.  If it avoids that conclusion by altering the object, standing packet, or endpoint role, it is carrier or domain shift.  Hence same-carrier endpoint-standing use without admissibility either does no endpoint work, exits the domain, or creates second endpoint governance.  K6 is the local exclusion of that split.
\end{proof}

\begin{lemma}[Local necessity of K11]\label{lem:k11-local-necessity}
Under \(C_0(M)\) plus minimal report evaluability, a theorem-bearing endpoint report that exceeds its carrier, standing, bridge, and support is unsupported endpoint authority, silent redescription, or no-carrier transfer.
\end{lemma}
\begin{proof}
Minimal report evaluability says only that a report remains attached to the fixed target and can be classified as theorem-bearing, support-level, bookkeeping, or out-of-domain.  K11 is the strengthening from that minimal evaluability to no-overreporting: a report may express no more endpoint force than its fixed target, carrier, standing packet, bridge position, and tensor support authorize.  If a negative report records obstruction evidence only, it is proof support.  If it reports endpoint failure on the same carrier, it governs endpoint status.  If that endpoint force is not supplied by a lawful endpoint role or bridge support, the report overstates its support.  The overstatement is unsupported endpoint authority when it remains on the carrier, silent redescription when a surrogate is redescribed as endpoint readout, or no-carrier transfer when status is imported from another object.  K11 is therefore the local report-preservation condition required for non-degenerate same-carrier endpoint use.
\end{proof}

\begin{lemma}[Local necessity of K13]\label{lem:k13-local-necessity}
K13 does not deny epistemic uncertainty, open proof status, incomplete information, or review status.  It states that unknown, deferred, unsupported, silent, bookkeeping, and hidden-fifth statuses are report or epistemic statuses, not endpoint truth branches on a fixed theorem-bearing slot.
\end{lemma}
\begin{proof}
Target-slot fixation by itself does not include endpoint exhaustion.  It fixes what is being evaluated.  K13 supplies the additional fail-closed theorem: once the carrier, standing packet, endpoint role, and admissibility boundary are fixed, a theorem-bearing endpoint branch is admitted for endpoint standing or not.  Unknown or deferred status may describe the state of a proof attempt; unsupported status may describe lack of tensor support; bookkeeping may organize notation; review status may describe reception.  None of these statuses occupies the endpoint truth slot.  Treating any of them as a third endpoint truth branch leaves the endpoint no longer determinate under same-carrier theorem use.  K13 is the local fail-closed endpoint-exhaustion condition that prevents that degeneration.
\end{proof}

\begin{longtable}{L{0.18\linewidth}L{0.36\linewidth}L{0.36\linewidth}}
\toprule
\textbf{Node} & \textbf{Weakening attempt} & \textbf{Degeneracy cost}\\
\midrule
K5 & Permit two admissible endpoint interiors for the same sphere slot, one positive and one negative. & If they agree, the second is bookkeeping. If they disagree, the same endpoint has two governing interiors on the same carrier, producing a second classifier or domain split.\\
K6 & Permit negative endpoint standing without admissibility. & Standing becomes an endpoint authority outside the admissibility gate. If endpoint-resolving, it is hidden classifier governance.\\
K11 & Permit a negative report to exceed carrier, bridge, and support. & The report becomes unsupported endpoint authority, silent redescription, or no-carrier transfer.\\
K13 & Permit unknown, deferred, unsupported, or independently negative governance as a third endpoint status. & The endpoint slot is no longer determinate; the branch is admitted, rejected, support-only, carrier shift, bookkeeping, or second endpoint governance.\\
\bottomrule
\end{longtable}

\begin{theorem}[No strict same-carrier weakening permits independent negative governance]\label{thm:no-strict-weakening}
Under the core endpoint adequacy packet \(C_0(M)\), there is no negative-governance-permissive weakening \(W\) that preserves non-degenerate same-carrier Poincare endpoint use.
\end{theorem}
\begin{proof}
Assume such a weakening \(W\) exists, and let \(G_-(M)\) be the negative endpoint-governance act it permits.  Since \(G_-(M)\) is endpoint-resolving, it changes the endpoint status of the fixed sphere-readout slot.  Since it is not proof support or bookkeeping, it is theorem-bearing.  Since it is not carrier shift, it acts on the same \(M\).  Since it is not bridge completion, it does not occupy the positive sphere bridge.  Since it is not a lawful coequal target role, it cannot be a separately declared alternative endpoint.

There are then only four endpoint-relevant possibilities.  First, \(G_-(M)\) shares the same admissible interior as the sphere bridge.  Then it collapses into that bridge and cannot supply negative exclusion.  Second, it has a distinct admissible interior.  By Lemma \ref{lem:k5-local-necessity}, this creates plural endpoint governance, a second classifier, or a domain split.  Third, it has endpoint standing without admissibility.  By Lemma \ref{lem:k6-local-necessity}, this is support-only, carrier shift, or a second gate.  Fourth, it is neither admitted nor rejected as endpoint-bearing while being preserved as endpoint-resolving.  By Lemma \ref{lem:k13-local-necessity}, that is a third-status degeneration, not determinate endpoint use.  If the negative report exceeds carrier and bridge support at any stage, Lemma \ref{lem:k11-local-necessity} gives unsupported endpoint authority, silent redescription, or no-carrier transfer.

Thus any weakening that permits \(G_-(M)\) either reduces it to bridge completion, proof support, bookkeeping, carrier shift, or lawful coequal target role, or else reproduces the forbidden independent endpoint-governance structure.  In none of these cases does \(W\) preserve non-degenerate same-carrier Poincare endpoint use while permitting the negative governance required to refute the endpoint.  Therefore no such strict weakening survives.
\end{proof}

\begin{corollary}[Strict weakening either preserves the packet extensionally or exits endpoint use]\label{cor:strict-weakening}
Let \(W\) be a weaker regime preserving core endpoint adequacy.  If \(W\) does not permit independent endpoint governance, then it is extensionally equivalent to the local packet for endpoint purposes.  If \(W\) permits independent endpoint governance, then it fails reference, standing, admissibility, irreversibility, target-slot fixation, minimal report evaluability, no-overreporting, endpoint exhaustion, or same-carrier fidelity.
\end{corollary}
\begin{proof}
If \(W\) does not alter endpoint-bearing statuses, its differences are bookkeeping, support-level, or presentation-level for the endpoint and it is extensionally equivalent to the local packet at the endpoint slot.  If \(W\) does alter endpoint-bearing status by permitting independent negative governance, Theorem \ref{thm:no-strict-weakening} shows that it exits the non-degenerate same-carrier endpoint regime by one of the listed failures.
\end{proof}

\begin{audit}[Right-mode weakening challenge]
A critic may examine K5, K6, K11, and K13 node by node.  The successful challenge must define a weakening \(W\) that preserves fixed carrier, theorem-bearing target-slot fixation, minimal report evaluability, standing--admissibility discipline, and no carrier shift, while also showing how the negative branch avoids proof support, bridge completion, lawful coequal target role, bookkeeping, and second endpoint authority.  The weakening-resistance theorem proves that no such \(W\) survives without importing no-overreporting or endpoint-exhaustion failure.
\end{audit}

\section{Poincare Carrier Adequacy and Role Separation}\label{sec:carrier-adequacy}

\begin{definition}[Standard Poincare carrier instantiation]
\(\mathrm{StdPCCarrierInst}(M)\) means that the ordinary topological data used in the Poincare problem instantiate \(\CarrierPC(M)\): the same closed connected three-manifold, the same homeomorphism relation, the same simple-connectedness standing packet, and the same sphere-readout predicate.
\end{definition}

\begin{definition}[Poincare carrier adequacy]
\(\mathrm{PCCarrierAdeq}(M)\) means:
\begin{enumerate}[label=(\alph*)]
\item \(M\) is the same closed connected topological 3-manifold under evaluation;
\item \(\Homeo\) is the lawful redescription relation;
\item \(\pi_1(M)=0\) is a fixed homeomorphism-invariant standing condition;
\item \(\SphereRead(M)\) is the endpoint readout, not standing;
\item category, geometry, triangulation, decomposition, and surgery data are auxiliary unless lawfully bridged back to \(M\).
\end{enumerate}
\end{definition}

\begin{theorem}[Standard Poincare carrier instantiates the adequate carrier]\label{thm:std-carrier-adequate}
For every closed connected 3-manifold \(M\),
\[
\mathrm{StdPCCarrierInst}(M)\Rightarrow \PCCarrierInst(M)\wedge \mathrm{PCCarrierAdeq}(M).
\]
\end{theorem}
\begin{proof}
The standard Poincare problem fixes precisely a closed connected 3-manifold, the homeomorphism relation, the simple-connectedness hypothesis, and the endpoint predicate \(M\Homeo\Sthree\). Model choices, triangulations, smooth presentations, metrics, or decompositions may change representatives or proof routes, but they do not change the fixed topological endpoint unless a lawful carrier bridge is supplied. Hence the standard problem instantiates the adequate carrier.
\end{proof}

\begin{lemma}[Carrier is not standing]
The carrier \(\CarrierPC(M)\) is not identical to \(\StandPC(M)\).
\end{lemma}
\begin{proof}
The carrier fixes the object class: a closed connected 3-manifold under homeomorphism. The standing packet adds \(\pi_1(M)=0\). A closed connected 3-manifold need not be simply connected. Thus carrier and standing are distinct roles.
\end{proof}

\begin{lemma}[Standing is not readout]
\(\StandPC(M)\) is not identical to \(\SphereRead(M)\).
\end{lemma}
\begin{proof}
If simple-connectedness standing were definitionally identical to sphere-readout, Poincare would be a definition. The endpoint is theorem-level only because standing and readout are distinct.
\end{proof}

\begin{lemma}[Negative branch is not ignorance]
The report \(M\not\Homeo\Sthree\) is not a neutral absence of proof. In the fixed Poincare endpoint it is the native negative branch when standing is fixed.
\end{lemma}
\begin{proof}
A statement that \(M\) is not homeomorphic to \(\Sthree\) denies the target readout on the fixed endpoint slot. If used theorem-bearingly, it resolves the endpoint negatively. If not used theorem-bearingly, it is proof support, bookkeeping, or lower-status data. It is not mere silence.
\end{proof}

\begin{theorem}[Poincare role separation]
Carrier, standing, sphere-readout, bridge slot, positive bridge completion, and native negative endpoint branch are distinct roles in the AASC endpoint architecture.
\end{theorem}
\begin{proof}
Carrier fixes \(M\); standing adds simple connectedness; readout asserts homeomorphism to \(\Sthree\); the bridge slot is the typed transition from standing to readout; bridge completion is the endpoint result; and the native negative branch denies bridge completion under the same standing packet. Collapsing any two of these changes the proof burden.
\end{proof}

\section{Sphere Bridge Packet and Negative Branch Correspondence}\label{sec:negative-correspondence}

This section avoids any direct inference from non-homeomorphism to an admissible discriminator. It first places the ordinary counterexample branch into native normal form and then identifies that branch with bridge-image exclusion.

\begin{definition}[Poincare negative branch]
\[
\PCNeg(M)\Longleftrightarrow \StandPC(M)\wedge M\not\Homeo\Sthree.
\]
This is the ordinary pointwise counterexample branch. It contains no AASC status vocabulary.
\end{definition}

\begin{definition}[Standard Poincare negative normal form]
\[
\StdPCNeg(M)\Longleftrightarrow \ClosedThree(M)\wedge\Conn(M)\wedge\dim M=3\wedge\pi_1(M)=0\wedge M\not\Homeo\Sthree.
\]
\end{definition}

\begin{theorem}[Negative normal form correspondence]\label{thm:pcneg-normal-form}
Under \(\mathrm{PCCarrierAdeq}(M)\),
\[
\PCNeg(M)\Longleftrightarrow \StdPCNeg(M).
\]
\end{theorem}
\begin{proof}
By definition \(\StandPC(M)\) is the conjunction of closedness, connectedness, dimension three, and simple connectedness. Substituting the definition of \(\StandPC(M)\) into \(\PCNeg(M)\) yields \(\StdPCNeg(M)\), and conversely. No endpoint-governance conclusion is used.
\end{proof}

\begin{definition}[Sphere-bridge slot]
For a fixed carrier \(M\), \(\BridgeSlotPC(M)\) is the declared transition position from \(\StandPC(M)\) to \(\SphereRead(M)\) on the same carrier. The slot fixes the typed transition. It does not assert spherehood.
\end{definition}

\begin{definition}[Poincare sphere-bridge object]\label{def:pc-bridge-object}
For a fixed pointwise Poincare carrier, \(\BridgeObjPC(M)\) is the sphere-bridge object recording a lawful bridge from \(\StandPC(M)\) to \(\SphereRead(M)\). It consists of the following preservation and report-discipline clauses.
\begin{enumerate}[label=\textup{(B\arabic*)}]
\item \textbf{Same-manifold preservation.} The bridge preserves the same manifold \(M\), not a surgery product, flow-stage object, model manifold, decomposition piece, triangulation artifact, cover, completion, or auxiliary carrier unless a lawful carrier bridge returns to \(M\).
\item \textbf{Carrier preservation.} The bridge preserves closedness, connectedness, dimension three, the homeomorphism relation, and the fixed topological carrier \(\CarrierPC(M)\).
\item \textbf{Standing preservation.} The bridge preserves the same standing packet
\[
\ClosedThree(M)\wedge\Conn(M)\wedge\dim M=3\wedge\pi_1(M)=0.
\]
\item \textbf{Readout preservation.} The bridge preserves the declared readout slot
\[
\SphereRead(M)\Longleftrightarrow M\Homeo\Sthree.
\]
It does not replace sphere-readout by homology, curvature, Ricci-flow extinction, triangulation normal form, recognition output, or category-specific surrogate unless that surrogate is bridged back to \(\SphereRead(M)\).
\item \textbf{Category preservation.} Smooth, PL, and topological presentations are allowed as bridge or support material only when they preserve the same topological endpoint carrier and do not create a separate endpoint branch.
\item \textbf{UEAP report preservation.} The bridge does not report standing as readout, proof support as bridge completion, auxiliary geometric data as topological endpoint readout, or non-spherehood as endpoint defeat without the corresponding bridge or endpoint-use role.
\item \textbf{ATS anchor/tensor/skin preservation.} The bridge preserves the anchor data \((M,\Homeo,\StandPC,\SphereRead)\), treats sphere-readout and bridge occupation as tensor content, and prevents notation, triangulation, metric choice, handlebody presentation, decomposition language, or Ricci-flow stage labels from carrying endpoint force as skin.
\end{enumerate}
\end{definition}

\begin{proposition}[Sphere-bridge object is non-circular]\label{prop:bridge-object-noncircular}
The bridge object \(\BridgeObjPC(M)\) does not contain \(\SphereRead(M)\) by definition. It fixes the same-manifold, carrier, standing, readout-slot, category, UEAP, and ATS constraints under which bridge completion may be reported. Sphere-readout appears only when the bridge is completed or when all endpoint-resolving negative occupations are excluded.
\end{proposition}
\begin{proof}
If \(\BridgeObjPC(M)\) contained \(\SphereRead(M)\), the Poincare endpoint would be definitional at the bridge-object layer. The clauses of \(\BridgeObjPC(M)\) instead fix the admissible route conditions: same carrier, same standing packet, same readout slot, category discipline, report preservation, and anchor/tensor/skin separation. None of these clauses asserts \(M\Homeo\Sthree\). Completion is a later theorem-bearing bridge status, not part of the object definition.
\end{proof}

\begin{definition}[Bridge-complete sphere object]
\(\BridgeObjCompPC(M)\) means that \(\BridgeObjPC(M)\) is discharged by theorem-bearing same-carrier bridge completion and therefore reports \(\SphereRead(M)\) on the fixed carrier.
\end{definition}

\begin{theorem}[Bridge-complete sphere object yields sphere-readout]\label{thm:bridge-complete-sphere-readout}
\[
\BridgeObjCompPC(M)\Rightarrow \SphereRead(M).
\]
\end{theorem}
\begin{proof}
Bridge-complete status is, by definition, completion of the standing-to-readout bridge on the same fixed carrier. Its output is precisely sphere-readout for \(M\).
\end{proof}

\begin{theorem}[Fixed-carrier bridge failure is sphere-bridge image exclusion]\label{thm:fixed-bridge-failure-exclusion}
Under fixed standing,
\[
\StandPC(M)\wedge\neg\SphereRead(M)\Rightarrow \SphereBridgeImgExcl(M).
\]
\end{theorem}
\begin{proof}
By the definition of sphere-bridge image exclusion below, sphere-bridge image exclusion is fixed standing together with denial of the sphere-readout slot. Thus fixed-standing non-spherehood is exactly bridge-image exclusion. This theorem makes the local projection explicit; it adds no endpoint-governance conclusion.
\end{proof}

\begin{definition}[Sphere-bridge admissibility]
\(B^{\mathrm{adm}}_{PC}(M)\) holds when a proposed transition preserves the same \(M\), the same homeomorphism relation, the same standing packet, and the same readout slot, while excluding surrogate substitution, carrier transfer, post-selection repair, selector import, and second-classifier governance.
\end{definition}

\begin{definition}[Sphere-bridge completion]
\(\BridgeCompPC(M)\) holds when the bridge slot and admissibility layer are discharged by theorem-bearing traces supporting \(\SphereRead(M)\) on the fixed carrier.
\end{definition}

\begin{proposition}[Non-circular bridge packet]\label{prop:bridge-noncircular}
The bridge slot, bridge object, and admissibility layer do not contain spherehood by definition. Spherehood appears only when bridge completion is certified or when endpoint-resolving negative occupation is excluded.
\end{proposition}
\begin{proof}
If the slot or \(\BridgeObjPC(M)\) contained \(\SphereRead(M)\), Poincare would be definitional. If admissibility contained \(M\Homeo\Sthree\), no incomplete bridge route could be audited. The slot fixes the transition type, the bridge object fixes preservation and report discipline, admissibility fixes same-carrier discipline, and completion is the endpoint result.
\end{proof}

\begin{definition}[Sphere-bridge image exclusion]
\(\SphereBridgeImgExcl(M)\) means that the fixed Poincare standing packet is active, the sphere-readout bridge slot is fixed, but standing and sphere-readout fail common endpoint occupation. In native notation:
\[
\SphereBridgeImgExcl(M)\Longleftrightarrow \StandPC(M)\wedge\neg\SphereRead(M).
\]
The concept is bridge-facing: it is not merely two symbols. It says that the fixed standing packet and declared readout do not occupy a common bridge endpoint.
\end{definition}

\begin{theorem}[Standard negative branch as bridge exclusion]\label{thm:stdneg-bridge-excl}
Under \(\mathrm{PCCarrierAdeq}(M)\),
\[
\StdPCNeg(M)\Longleftrightarrow \SphereBridgeImgExcl(M).
\]
\end{theorem}
\begin{proof}
The standard negative normal form fixes the Poincare standing packet and denies sphere-readout. Since \(\SphereRead(M)\) is precisely \(M\Homeo\Sthree\), this is exactly the claim that the fixed standing packet and the sphere-readout slot fail common bridge occupation. Conversely, sphere-bridge image exclusion fixes standing and denies sphere-readout, which is the standard negative normal form. No AASC discriminator conclusion is used.
\end{proof}

\begin{theorem}[Poincare negative branch is bridge-image exclusion]\label{thm:pcneg-bridge-excl}
Under \(\mathrm{PCCarrierAdeq}(M)\),
\[
\PCNeg(M)\Longleftrightarrow \SphereBridgeImgExcl(M).
\]
\end{theorem}
\begin{proof}
Combine Theorem \ref{thm:pcneg-normal-form} with Theorem \ref{thm:stdneg-bridge-excl}.
\end{proof}

\begin{definition}[Theorem-level sphere-status discriminator]
\(\ThmSphereDisc(M)\) means a theorem-level classifier assigning endpoint status to the fixed standing/readout pair by declaring either sphere-bridge completion or sphere-bridge image exclusion. It is not a preferred triangulation, Ricci-flow stage, metric selector, decomposition artifact, homology surrogate, recognition algorithm, or mere restatement of notation. It classifies endpoint status of already-fixed roles.
\end{definition}

\begin{theorem}[Endpoint-bearing bridge exclusion induces theorem-level sphere-status discrimination]\label{thm:bridge-excl-disc}
\[
\SphereBridgeImgExcl(M)\wedge\OfficialPCNegativeResolution(M)\Rightarrow \ThmSphereDisc(M).
\]
\end{theorem}
\begin{proof}
Sphere-bridge image exclusion as a descriptive observation is not automatically endpoint-status discrimination. The additional hypothesis that it is used as the official negative pointwise resolution makes it theorem-bearing endpoint content. It is not standing alone, because standing is only \(\StandPC(M)\). It is not sphere-readout, because it denies \(\SphereRead(M)\). It is not bridge completion, because it declares common bridge occupation absent. Therefore it classifies the status of the fixed standing/readout pair as bridge-excluded. That is theorem-level sphere-status discrimination.
\end{proof}

\section{ATS and UEAP Audit of Negative Endpoint Use}\label{sec:ats-ueap-negative}

\begin{definition}[ATS Poincare layers]
For the Poincare endpoint, the anchor layer fixes \(M\), \(\Homeo\), closedness, connectedness, dimension three, simple connectedness, \(\SphereRead\), and the bridge slot. The tensor layer contains load-bearing endpoint structures: standing, sphere-readout, bridge completion, bridge-image exclusion, theorem-level sphere-status discrimination, endpoint-status governance, route-locus use, and no-second-classifier closure. The skin layer contains notation, diagrams, representative choices, preferred triangulations, numerical experiments, heuristic reports, status labels, and non-endpoint descriptive reformulations.
\end{definition}

\begin{theorem}[ATS classification of negative branch use]\label{thm:ats-negative-classification}
On the fixed Poincare carrier, a \(\PCNeg\)-like statement has only three ATS statuses: skin/support if not endpoint-resolving, tensor if it changes Poincare endpoint status, or carrier shift if the fixed carrier is not preserved. No fourth ATS status exists.
\end{theorem}
\begin{proof}
If the statement does not alter the endpoint report, it is proof support, descriptive data, or skin relative to the endpoint. If it changes the endpoint report on the fixed carrier, it performs load-bearing theorem work and is tensor. If it obtains its effect by changing the manifold, equivalence relation, standing packet, category, compactness status, or endpoint predicate, it is carrier shift rather than same-endpoint use. These cases exhaust the ATS role alternatives.
\end{proof}

\begin{definition}[UEAP negative report bound]
A Poincare report may not exceed its target, carrier, standing, and bridge support. Thus \(\StandPC(M)\) reports standing only; \(\SphereRead(M)\) reports readout; bridge completion reports common occupation; and a negative branch offered as official endpoint resolution must have a tensor-level endpoint role.
\end{definition}

\begin{theorem}[UEAP bound for official negative Poincare reports]\label{thm:ueap-negative-bound}
An official negative Poincare report \(M\not\Homeo\Sthree\) exceeds carrier standing alone and readout alone. Therefore, on the fixed carrier, it must be one of: proof support only, endpoint-status tensor governance, carrier shift, bookkeeping, or the hidden fifth case of endpoint resolution without governance.
\end{theorem}
\begin{proof}
Standing supplies \(\StandPC(M)\); readout supplies \(\SphereRead(M)\). The negative endpoint report compares the fixed standing packet with the sphere-readout slot and declares common bridge occupation absent. This exceeds either component role taken alone. UEAP therefore requires a lawful report status for the stronger act. If the act is not endpoint-resolving, it is support or bookkeeping. If it changes target, it is carrier shift. If it resolves the endpoint while claiming no endpoint-status work, it is the hidden fifth case. Otherwise it is tensor-level endpoint-status governance.
\end{proof}

\begin{audit}[Ordinary negative data versus endpoint-governing negative use]
The manuscript does not forbid descriptive topology or counterexample-like expressions. It classifies the role such data take when they are used to resolve the official Poincare endpoint.
\begin{longtable}{L{0.48\linewidth}L{0.42\linewidth}}
\toprule
\textbf{Claimed negative use} & \textbf{Classification}\\
\midrule
Local evidence, hypothetical obstruction, computation, decomposition fact, or proof support & Legitimate support; not endpoint-resolving.\\
A theorem proving \(M\not\Homeo\Sthree\) for fixed \(M\) under official pointwise endpoint audit & Endpoint-status governance.\\
Endpoint-resolving negative theorem that denies governance & Hidden fifth case; impossible.\\
Negative conclusion obtained by changing \(M\), category, compactness, boundary status, or equivalence relation & Carrier/domain shift.\\
Expression with no endpoint effect & Bookkeeping.\\
\bottomrule
\end{longtable}
\end{audit}

\section{Negative-Use Classification and No Hidden Fifth Case}\label{sec:no-fifth}

\begin{definition}[Poincare negative use kind]
A \(\PCNeg\)-like statement has one of five use kinds:
\begin{enumerate}[label=(\roman*)]
\item proof-support observation: data or local evidence used inside a proof but not offered as official Poincare endpoint resolution;
\item endpoint-resolving negative theorem: non-spherehood is offered as the official negative resolution of the Poincare endpoint on the fixed carrier;
\item endpoint-resolving non-governance: non-spherehood allegedly resolves the endpoint while doing no endpoint-status governance;
\item carrier-changing negative claim: non-spherehood is obtained only after changing \(M\), category, compactness, boundary status, equivalence relation, or endpoint role;
\item bookkeeping-only expression: the expression is recorded but does no endpoint work.
\end{enumerate}
\end{definition}

\begin{definition}[Negative-use classification]
The use-kind classification is:
\begin{longtable}{L{0.43\linewidth}L{0.47\linewidth}}
\toprule
\textbf{Use kind} & \textbf{Classification}\\
\midrule
Proof-support observation & Legitimate as support; not endpoint-resolving.\\
Endpoint-resolving negative theorem & Tensor endpoint-status governance.\\
Endpoint-resolving non-governance & Hidden fifth case; impossible.\\
Carrier-changing negative claim & Domain shift; not same endpoint.\\
Bookkeeping-only expression & No endpoint occupation.\\
\bottomrule
\end{longtable}
\end{definition}

\begin{theorem}[Endpoint-resolving non-governance is impossible]\label{thm:no-hidden-fifth}
There is no endpoint-resolving, fixed-carrier, non-governing Poincare negative use.
\end{theorem}
\begin{proof}
By Theorem \ref{thm:ats-negative-classification}, an endpoint-resolving same-carrier statement is tensor. By Theorem \ref{thm:ueap-negative-bound}, an official negative report is proof support, endpoint-status tensor governance, carrier shift, bookkeeping, or the hidden fifth case. Proof support and bookkeeping do not resolve the endpoint; carrier shift is not same-endpoint use. Thus a statement that resolves the endpoint and refuses endpoint-status governance claims an impossible role: tensor by endpoint effect and non-tensor by self-description. No such use survives.
\end{proof}

\begin{definition}[Official Poincare negative resolution]
\(\OfficialPCNegativeResolution(M)\) means theorem-bearing use of \(\PCNeg(M)\) as the official negative endpoint resolution of the pointwise Poincare assertion on the fixed carrier.
\end{definition}

\begin{theorem}[Pointwise negative endpoint branch]\label{thm:pcneg-official-negative}
Under the official endpoint target under audit,
\[
\begin{aligned}
&\PCEndpoint\wedge\PCCarrierInst(M)\wedge\PCNeg(M)\\
&\qquad\Rightarrow\OfficialPCNegativeResolution(M).
\end{aligned}
\]
\end{theorem}
\begin{proof}
The official pointwise endpoint for the fixed carrier is \(\SphereRead(M)\). Under carrier adequacy and Theorem \ref{thm:pcneg-bridge-excl}, \(\PCNeg(M)\) is exactly the fixed-carrier complement branch: standing is active and sphere-readout is denied. If that branch is theorem-bearing in the official endpoint target under audit, it resolves the pointwise endpoint negatively. If it is not theorem-bearing, it is proof support, bookkeeping, or incomplete expression. If it changes carrier, it is not the same pointwise endpoint. Therefore, under official endpoint audit and fixed carrier instantiation, \(\PCNeg(M)\) is the official negative resolution.
\end{proof}

\begin{definition}[Poincare endpoint-status governance]
\(\EndpointGov(M)\) means tensor-level endpoint-status governance of the Poincare sphere bridge by a theorem-level negative-status discriminator that is not itself positive sphere-bridge completion.
\end{definition}

\begin{theorem}[Official negative resolution is endpoint-status governance]\label{thm:official-negative-governance}
\[
\OfficialPCNegativeResolution(M)\Rightarrow\EndpointGov(M).
\]
\end{theorem}
\begin{proof}
Official negative resolution settles the Poincare endpoint negatively. By Theorem \ref{thm:no-hidden-fifth}, endpoint-resolving non-governance is impossible. It is not proof support only, because it resolves the endpoint. It is not bookkeeping, because bookkeeping has no endpoint effect. It is not carrier shift, because official negative resolution is on the fixed carrier. It is not sphere-bridge completion, because it denies bridge occupation. Hence the only remaining same-carrier endpoint-resolving role is endpoint-status governance.
\end{proof}


\section{Non-Explosiveness of Negative-Branch Exclusion}\label{sec:non-explosiveness}

The exclusion proved here is not a global ban on negative mathematics. It is a narrow fixed-endpoint constraint. AASC permits negative results when the negative branch occupies a lawful negative target role, functions as proof support, diagnoses carrier shift, supplies an obstruction inside a bridge route, or belongs to an explicitly fixed coequal bivalent endpoint. What is excluded in the Poincare proof is more specific: a same-carrier endpoint-resolving negative branch that does not occupy the positive sphere bridge, does not remain proof support or bookkeeping, does not change carrier, and nevertheless governs the fixed sphere-readout endpoint from outside bridge occupation.

\begin{theorem}[Negative results are not globally forbidden]\label{thm:negative-not-globally-forbidden}
AASC does not forbid negative mathematical results as such. It forbids, on a fixed endpoint domain, an endpoint-resolving negative use that simultaneously preserves the same carrier, denies bridge occupation, refuses proof-support or bookkeeping status, and supplies endpoint-status authority independently of the admitted bridge.
\end{theorem}
\begin{proof}
A negative result attached to a lawful negative target role is governed by that role. A negative obstruction used inside a proof route is support until it resolves the endpoint. A negative statement obtained by changing carrier is a domain shift. A negative datum with no endpoint effect is bookkeeping. None of these cases is excluded by no-independent-discriminator closure. The excluded case is the residue after those legitimate roles are removed: same-carrier endpoint resolution by a negative branch that does not occupy the positive bridge but still governs endpoint status. That residue is independent same-domain endpoint authority, which the kernel excludes.
\end{proof}

\begin{remark}[No arbitrary theorem explosion]\label{rem:no-explosion}
The proof does not proceed by naming an arbitrary universal target and banning its complement. The no-independent-discriminator theorem applies only after the fixed carrier, standing packet, endpoint role, admissibility boundary, bridge-role typing, no-hidden-fifth classification, and endpoint-forced kernel are all in force. A false target with a lawful coequal negative endpoint role, or a target whose negative branch remains ordinary proof support or carrier diagnosis, does not instantiate the forbidden Poincare structure.
\end{remark}

\section{Independent Sphere Discriminator and Its Exclusion}\label{sec:independent-disc}

\begin{definition}[Independent Poincare sphere discriminator]
\(\IndependentSphereDisc(M)\) means a theorem-level sphere-status discriminator that:
\begin{enumerate}[label=(\alph*)]
\item preserves the same \(M\), closedness, connectedness, dimension three, simple-connectedness standing, and homeomorphism endpoint role;
\item does not occupy positive sphere-bridge completion;
\item nevertheless governs whether the sphere endpoint is occupied or excluded;
\item is not a lawful carrier reset, proof-support observation, bookkeeping act, selector output, surrogate readout, or bridge-complete certificate.
\end{enumerate}
\end{definition}

\begin{theorem}[Endpoint-status governance induces independent sphere discrimination]\label{thm:governance-independent-disc}
\[
\EndpointGov(M)\Rightarrow\IndependentSphereDisc(M).
\]
\end{theorem}
\begin{proof}
Endpoint-status governance is independent precisely when it settles the fixed sphere endpoint while failing to occupy the positive bridge and while preserving the same carrier. The possible legitimate roles are exhausted by the following audit.
\begin{longtable}{L{0.35\linewidth}L{0.55\linewidth}}
\toprule
\textbf{Possible role} & \textbf{Why endpoint-resolving \(\PCNeg\) is not that role}\\
\midrule
Carrier standing & It reports only closed connected simply connected 3-manifold status; it does not by itself settle sphere-readout.\\
Sphere-readout & It denies \(\SphereRead(M)\).\\
Sphere-bridge completion & It declares common bridge occupation absent.\\
Bridge support & Support may contribute to a proof, but support alone does not resolve the endpoint.\\
Carrier shift & Official endpoint use fixes \(\CarrierPC(M)\), so carrier-changing comparison is not same endpoint use.\\
Bookkeeping & Bookkeeping has no endpoint effect and therefore cannot resolve the endpoint.\\
Ordinary counterexample without governance & This is the hidden fifth case, excluded by Theorem \ref{thm:no-hidden-fifth}.\\
Lawful bridge certificate & A bridge certificate would occupy the positive sphere bridge and yield sphere-readout, not bridge exclusion.\\
\bottomrule
\end{longtable}
After these roles are removed, a same-carrier endpoint-resolving negative theorem governs the endpoint while not occupying the positive bridge. That is exactly an independent Poincare sphere discriminator.
\end{proof}

\begin{remark}[Scope of the no-independent-discriminator theorem]\label{rem:no-independent-scope}
The theorem below does not say that negative mathematical results are impossible. A negative theorem is not independent endpoint governance merely because it is negative. It becomes independent endpoint governance only when it is used to resolve the same fixed endpoint while not occupying a lawful negative target role, not remaining proof support, not changing carrier, and not supplying bridge completion. On the fixed Poincare endpoint carrier, that is the forbidden structure.
\end{remark}

\begin{theorem}[No independent same-domain Poincare sphere discriminator]\label{thm:no-independent-sphere-disc}
For every fixed official Poincare endpoint carrier,
\[
\KPC(M)\Rightarrow\neg\IndependentSphereDisc(M).
\]
\end{theorem}
\begin{proof}
An independent sphere discriminator would govern the same fixed sphere endpoint while not occupying the positive sphere bridge and not changing carrier. It would therefore supply a second same-domain endpoint-status authority over the standing-to-readout transition. Lemma \ref{lem:plural-routes-not-interiors} permits ordinary plurality of proof routes and support data; what is barred is plural endpoint-governing authority for the same fixed endpoint slot. By K5, plural admissible interiors for that slot are excluded. By K6, endpoint standing cannot float outside admissibility. By K11, reports may not exceed carrier and bridge support. By K13, endpoint-bearing branches are exhausted once carrier, standing, endpoint role, and admissibility boundary are fixed. The weakening-resistance theorem, Theorem \ref{thm:no-strict-weakening}, also shows that a purported weaker same-carrier regime permitting such governance is not a non-degenerate endpoint regime. Thus the alleged discriminator has only five possible statuses: positive bridge occupation, proof support only, carrier/domain reset, bookkeeping, or forbidden second-classifier governance. The definition of \(\IndependentSphereDisc\) excludes the first four. The fifth is inadmissible by K5, K6, K11, and K13. Therefore no independent same-domain Poincare sphere discriminator survives.
\end{proof}

\begin{corollary}[Native negative branch is excluded under official endpoint use]\label{cor:pcneg-excluded}
For every \(M\),
\[
\OfficialPCEndpointUse(M)\Rightarrow\neg\PCNeg(M).
\]
\end{corollary}
\begin{proof}
Assume official pointwise endpoint use and the native negative branch for \(M\). Proposition \ref{prop:endpoint-use-carrier} gives carrier instantiation, and Theorem \ref{thm:endpoint-use-kernel} gives \(\KPC(M)\). Theorem \ref{thm:pcneg-official-negative} gives official negative resolution. Theorem \ref{thm:official-negative-governance} gives endpoint-status governance, and Theorem \ref{thm:governance-independent-disc} gives independent sphere discrimination. This contradicts Theorem \ref{thm:no-independent-sphere-disc}. Hence \(\neg\PCNeg(M)\).
\end{proof}


\section{Local Reductio Engine and Sphere-Readout Closeout}\label{sec:local-reductio}

The global negative-resolution route in the previous sections remains the route audit for official endpoint-defeating use. The endpoint closeout below uses the local exact-complement route: the temporary assumption \([\neg\SphereRead(M)]_i\) is opened only as the live countercase against sphere-readout and is not promoted to a global negative endpoint result.

\begin{definition}[Exact-complement non-spherehood annotation]
The symbol
\[
\ReductioCountercasePC(\neg\SphereRead(M);\SphereRead(M))
\]
denotes the proof-act annotation generated when the sole local assumption \([\neg\SphereRead(M)]_i\) is opened as the live exact-complement countercase against \(\SphereRead(M)\) under fixed \(\StandPC(M)\). In native notation, this is \([M\not\Homeo\Sthree]_i\). The annotation records proof role only; it is not an additional object-level premise.
\end{definition}

\begin{theorem}[Annotation is not a premise]\label{thm:pc-annotation-not-premise}
The exact-complement annotation does not add a second conjunct to the discharged assumption. The only object-level assumption opened and later discharged is \(\neg\SphereRead(M)\); the annotation records only that this assumption is being used as the live countercase against \(\SphereRead(M)\).
\end{theorem}
\begin{proof}
A proof-act annotation records why an assumption is being used in a subproof. It does not add a second proposition to the assumption list. The local subproof opens only \(\neg\SphereRead(M)\); the annotation records that this assumption is the exact complement of the target in the fixed pointwise endpoint proof.
\end{proof}

\begin{definition}[Local Poincare countercase]
\(\LocalCountercasePC(M)\) means that the assumption \([\neg\SphereRead(M)]_i\) is the live exact-complement countercase against \(\SphereRead(M)\) under fixed \(\StandPC(M)\).
\end{definition}

\begin{theorem}[Exact-complement non-spherehood has local countercase standing]\label{thm:local-countercase-standing}
Under \(\PCCarrierInst(M)\), \(\StandPC(M)\), and the exact-complement annotation, the local assumption \([\neg\SphereRead(M)]_i\) has local countercase standing: \(\LocalCountercasePC(M)\).
\end{theorem}
\begin{proof}
The fixed pointwise endpoint has one readout slot, \(\SphereRead(M)\). Its exact complement in the local proof is \(\neg\SphereRead(M)\). Opening that complement as the branch whose survival would block endpoint closure gives it temporary local countercase standing. This is not global theoremhood of the negative branch and not official negative endpoint resolution.
\end{proof}

\begin{definition}[Local endpoint counterforce]
The predicate \(\EndpointCounterforcePC(M)\) records that the live local countercase blocks the sphere-readout target on the fixed carrier.
\end{definition}

\begin{theorem}[Local countercase has endpoint counterforce]\label{thm:local-counterforce}
\[
\LocalCountercasePC(M)\Rightarrow \EndpointCounterforcePC(M).
\]
\end{theorem}
\begin{proof}
A live exact-complement countercase is introduced precisely to test whether the target can be blocked. That endpoint-facing role is endpoint counterforce inside the subproof. The conclusion is local and does not assert global official negative endpoint resolution.
\end{proof}

\begin{definition}[Local sphere-bridge image exclusion]
\(\SphereBridgeImgExclLoc(M)\) means that sphere-bridge image exclusion is used inside the exact-complement reductio subproof as the endpoint-defeating negative branch on the fixed Poincare carrier.
\end{definition}

\begin{theorem}[Local countercase projects to local bridge-image exclusion]\label{thm:local-countercase-bridge-exclusion}
\[
\StandPC(M)\wedge \neg\SphereRead(M)\wedge \EndpointCounterforcePC(M)
\Rightarrow \SphereBridgeImgExclLoc(M).
\]
\end{theorem}
\begin{proof}
Under fixed standing, \(\neg\SphereRead(M)\) is exactly failure of common occupation of the standing-to-readout bridge image. Endpoint counterforce says this bridge-image exclusion is being used as the local branch whose survival would block endpoint closure. That is precisely local sphere-bridge image exclusion.
\end{proof}

\begin{theorem}[Local bridge-image exclusion induces independent sphere discrimination]\label{thm:local-bridge-exclusion-independent}
\[
\SphereBridgeImgExclLoc(M)\Rightarrow \IndependentSphereDisc(M).
\]
\end{theorem}
\begin{proof}
Local bridge-image exclusion is endpoint-facing, fixed-carrier, non-bridge-complete sphere-status work. It is not carrier standing alone, not sphere-readout, not bridge completion, not proof support only, not bookkeeping, and not carrier shift. By the same role exhaustion used in Theorem \ref{thm:governance-independent-disc}, the remaining same-domain endpoint-status role is independent sphere discrimination.
\end{proof}

\begin{theorem}[Local no-independent closure]\label{thm:local-no-independent}
\[
\KPC(M)\Rightarrow \neg\IndependentSphereDisc(M).
\]
\end{theorem}
\begin{proof}
This is Theorem \ref{thm:no-independent-sphere-disc} applied to the fixed pointwise carrier. The no-independent closure does not distinguish global official resolution from local endpoint-facing countercase use when the same carrier, standing packet, readout slot, and endpoint role are preserved.
\end{proof}

\begin{corollary}[Local non-spherehood countercase contradiction]\label{cor:local-countercase-contradiction}
Under \(\KPC(M)\), \(\PCCarrierInst(M)\), and \(\StandPC(M)\), the local assumption \([\neg\SphereRead(M)]_i\) yields contradiction.
\end{corollary}
\begin{proof}
Open \([\neg\SphereRead(M)]_i\) as the exact-complement countercase against \(\SphereRead(M)\). The annotation gives \(\LocalCountercasePC(M)\) by Theorem \ref{thm:local-countercase-standing}. Theorem \ref{thm:local-counterforce} gives endpoint counterforce, and Theorem \ref{thm:local-countercase-bridge-exclusion} gives \(\SphereBridgeImgExclLoc(M)\). Theorem \ref{thm:local-bridge-exclusion-independent} gives \(\IndependentSphereDisc(M)\), while Theorem \ref{thm:local-no-independent} gives its negation. Contradiction follows.
\end{proof}

\begin{theorem}[Annotation discharge yields sphere-readout]\label{thm:annotation-discharge-sphere-readout}
If the annotated exact-complement subproof from \([\neg\SphereRead(M)]_i\) yields contradiction, then the discharged result is \(\SphereRead(M)\).
\end{theorem}
\begin{proof}
By Theorem \ref{thm:pc-annotation-not-premise}, the annotation is not an added object-level premise. The sole discharged object-level assumption is \(\neg\SphereRead(M)\). Classical reductio gives \(\neg\neg\SphereRead(M)\), hence \(\SphereRead(M)\).
\end{proof}


\section{Endpoint Closure}\label{sec:main}

\begin{theorem}[Local sphere-readout closeout]\label{thm:local-sphere-readout-closeout}
On the fixed official endpoint target,
\[
\PCEndpoint\wedge\StandPC(M)\Rightarrow \SphereRead(M).
\]
\end{theorem}
\begin{proof}
Assume \(\PCEndpoint\) and \(\StandPC(M)\). By Theorem \ref{thm:universal-binds-pointwise}, the pointwise carrier is under \(\OfficialPCEndpointUse(M)\) and \(\PCCarrierInst(M)\). Theorem \ref{thm:endpoint-use-kernel} gives \(\KPC(M)\). Open the exact-complement assumption \([\neg\SphereRead(M)]_i\). Corollary \ref{cor:local-countercase-contradiction} gives contradiction in the annotated subproof. Theorem \ref{thm:annotation-discharge-sphere-readout} discharges the sole object-level assumption and yields \(\SphereRead(M)\).
\end{proof}

\begin{corollary}[Global negative-branch exclusion remains available]\label{cor:global-negative-branch-exclusion}
The global route audit still excludes the native negative branch:
\[
\PCEndpoint\wedge\StandPC(M)\Rightarrow \neg\PCNeg(M).
\]
\end{corollary}
\begin{proof}
Assume \(\PCEndpoint\) and \(\StandPC(M)\). Theorem \ref{thm:local-sphere-readout-closeout} gives \(\SphereRead(M)\). Since \(\PCNeg(M)\) is \(\StandPC(M)\wedge\neg\SphereRead(M)\), \(\PCNeg(M)\) is impossible. This corollary records the global negative-branch audit; the endpoint closeout above uses local exact-complement discharge.
\end{proof}

\begin{theorem}[AASC Poincare closure theorem]\label{thm:aasc-poincare}
In AASC endpoint mode, target fixation plus the fixed standing packet force sphere-readout:
\[
\PCEndpoint\wedge\ClosedThree(M)\wedge\Conn(M)\wedge\dim M=3\wedge\pi_1(M)=0\Rightarrow M\Homeo\Sthree.
\]
\end{theorem}
\begin{proof}
Assume \(\PCEndpoint\) and let \(M\) be closed, connected, three-dimensional, and simply connected. Then \(\StandPC(M)\) holds. Theorem \ref{thm:local-sphere-readout-closeout} gives \(\SphereRead(M)\). By definition of sphere-readout, \(M\Homeo\Sthree\). The antecedent \(\PCEndpoint\) is the theorem-target fixation of this manuscript, not an assumption of \(\SphereRead(M)\).
\end{proof}

\begin{corollary}[Smooth category form]
\leavevmode\par
In AASC endpoint mode, under \(\PCEndpoint\), every closed, connected, simply connected smooth three-manifold is diffeomorphic to \(\Sthree\), subject to the standard three-dimensional category bridge.
\end{corollary}
\begin{proof}
Assume the same endpoint-target fixation \(\PCEndpoint\) used in Theorem \ref{thm:aasc-poincare}. Apply Theorem \ref{thm:aasc-poincare} to the underlying topological carrier. Dimension-three category discipline prevents smooth or PL presentation drift from generating a separate endpoint branch; the standard bridge supplies the smooth presentation after topological closure \cite{Moise1977}.
\end{proof}

\section{Route-Locus Exhaustion for Negative Endpoint Use}\label{sec:route-locus}

The route-locus material is used here only in its narrowed endpoint-governance role. It does not claim that every possible topological truth must first present itself as an independent discriminator. It classifies every same-carrier endpoint-resolving negative route.

\begin{definition}[Endpoint-resolving negative route]
An endpoint-resolving negative route for \(M\) is a proof trace, invariant, obstruction, construction, classifier, or report that is used to make \(\PCNeg(M)\) the official negative resolution of the pointwise Poincare endpoint on \(\CarrierPC(M)\).
\end{definition}

\begin{definition}[Poincare route loci]
The route-locus basis consists of:
\begin{enumerate}[label=L\arabic*.]
\item carrier identity;
\item fundamental group;
\item homology/cohomology;
\item local-to-global assembly;
\item decomposition/prime structure;
\item boundary/end/compactness;
\item category/smooth/PL/topological structure;
\item geometric or analytic auxiliary carrier;
\item modification/repair history;
\item representative selection;
\item endpoint-classifier or negative-status governance.
\end{enumerate}
\end{definition}

\begin{theorem}[Route-locus exhaustion]\label{thm:route-locus-exhaustion}
Every same-carrier endpoint-resolving negative route for the Poincare endpoint enters at least one of L1--L11. If it enters none, it does no endpoint work and cannot resolve the official endpoint.
\end{theorem}
\begin{proof}
A route that resolves the endpoint negatively must do some endpoint-status work. If it changes the object under evaluation, it enters carrier identity. If it changes loop standing, it enters the fundamental group locus. If it uses algebraic topology, it enters homology/cohomology. If it assembles local data into global obstruction, it enters local/global assembly. If it uses splitting, gluing, or prime data, it enters decomposition. If it invokes boundary, puncture, end, or noncompact behavior, it enters boundary/end/compactness. If it changes smooth/PL/topological presentation, it enters category. If it imports metric, curvature, Ricci flow, PDE, or geometric decomposition data, it enters auxiliary carrier. If it uses surgery, cutting, collapse, reinterpretation, or after-the-fact validation, it enters modification/repair history. If it decides by a preferred triangulation, metric, handlebody, presentation, or representative, it enters selector authority. If it changes the rule by which spherehood is decided, it enters endpoint-classifier or negative-status governance. If none of these occurs, the route preserves carrier, standing, algebraic content, local/global status, decomposition, boundary/end status, category, auxiliary carrier, modification history, representative choice, and endpoint classifier; it has no remaining endpoint work and cannot resolve the endpoint.
\end{proof}

\begin{longtable}{L{0.26\linewidth}L{0.38\linewidth}L{0.26\linewidth}}
\toprule
\textbf{Feature} & \textbf{Endpoint work attempted} & \textbf{Disposition}\\
\midrule
Fundamental group & Supply negative endpoint resolution by loop obstruction. & Blocked by \(\pi_1(M)=0\) or routes out of L2.\\
Homology/cohomology & Replace homeomorphism readout by algebraic surrogate. & Support only unless lawfully bridged; endpoint use routes to governance.\\
Local charts & Promote local manifold data to global non-sphere status. & Local data are carrier data; global obstruction routes elsewhere.\\
Decomposition & Use splitting, prime residue, or gluing structure as negative endpoint. & Carrier drift, group/homology route, or governance route.\\
Boundary/end & Distinguish by puncture, end, or noncompact escape. & Contradicts closed carrier or resets scope.\\
Category & Use smooth/PL/topological mismatch as endpoint obstruction. & Dimension-three category discipline blocks independent residue.\\
Geometry/Ricci flow & Carry endpoint status through auxiliary metric or flow carrier. & No carrier transfer without bridge; bridge completion is positive or route-locus data.\\
Surgery/repair & Validate a negative conclusion after modification. & Carrier drift or no post-selection repair.\\
Selector & Decide by preferred triangulation, handlebody, metric, or presentation. & AMetric no-selector boundary.\\
Classifier & Govern the endpoint by a rule beside the bridge. & No independent same-domain discriminator.\\
\bottomrule
\end{longtable}

\section{Route-Locus Collapse Audits}\label{sec:locus-collapses}

\subsection{L1: Carrier identity}
\begin{theorem}[Carrier-drift collapse]
Any endpoint-resolving negative route that changes the target from the fixed \(M\) to a related manifold, model, decomposition piece, geometric flow stage, triangulation artifact, or surgery product fails as same-domain negative endpoint evidence unless carrier preservation is certified.
\end{theorem}
\begin{proof}
The endpoint concerns \(M\). A report about another object is not a report about \(M\). If a route claims the other object carries status back to \(M\), the route must supply bridge data. Without such data it is carrier transfer; with such data the bridge, not the drifted object, is load-bearing and enters another route locus.
\end{proof}

\subsection{L2: Fundamental group}
\begin{theorem}[Fundamental-group collapse]
No fundamental-group negative route survives under \(\pi_1(M)=0\).
\end{theorem}
\begin{proof}
A route using nontrivial loop obstruction contradicts the standing hypothesis. A route not using loop obstruction is not located in L2. Hence L2 supplies no endpoint-resolving negative route on the fixed carrier.
\end{proof}

\subsection{L3: Homology/cohomology}
\begin{theorem}[Homology/cohomology route collapse]
Homology or cohomology may be legitimate proof support, but it cannot by itself resolve the Poincare endpoint negatively on the fixed carrier.
\end{theorem}
\begin{proof}
Homological data are algebraic summaries, not homeomorphism readout. Agreement with \(\Sthree\) gives no negative endpoint resolution. Disagreement either contradicts the fixed standing packet through standard algebraic-topological constraints or requires a further route explaining why the algebraic datum decides homeomorphism status. That further route enters another locus or endpoint governance. Thus homology/cohomology alone does not survive as same-carrier endpoint-resolving negative use.
\end{proof}

\subsection{L4: Local-to-global assembly}
\begin{theorem}[Local/global route collapse]
Local Euclidean structure cannot supply a negative endpoint resolution without a global compatibility obstruction, and any such obstruction routes to another locus.
\end{theorem}
\begin{proof}
Both \(M\) and \(\Sthree\) have local 3-ball charts. Local structure is carrier data. If the route alleges failure of global assembly as \(\Sthree\), the load-bearing datum is global compatibility, not locality; it routes to group, homology, decomposition, category, geometry, selector, carrier, or classifier governance. Therefore no independent local/global negative route survives.
\end{proof}

\subsection{L5: Decomposition}
\begin{theorem}[Decomposition route collapse]
A decomposition-based negative route either changes carrier, routes to another locus, or becomes endpoint-status governance.
\end{theorem}
\begin{proof}
A decomposition route must use splitting, gluing graph, prime residue, neck structure, or piece structure. If it changes the target object, L1 collapses it. If it introduces loop or gluing obstruction, it routes to L2 or L4. If it introduces algebraic residue, it routes to L3. If it decides the endpoint by rule rather than bridge, it routes to L11. If it does none of these, it performs no endpoint work.
\end{proof}

\subsection{L6: Boundary/end/compactness}
\begin{theorem}[Boundary/end route collapse]
Boundary, puncture, noncompactness, end structure, or escape-to-infinity data cannot resolve the closed Poincare carrier negatively without changing the carrier.
\end{theorem}
\begin{proof}
The carrier fixes \(M\) as closed and compact without boundary. Any boundary, puncture, or end datum contradicts the carrier or changes scope. Such a route is carrier drift, not same-endpoint negative resolution.
\end{proof}

\subsection{L7: Category and smooth structure}
\begin{theorem}[Category route collapse]
Smooth, PL, or topological category variation in dimension 3 does not supply independent same-domain negative endpoint governance once the topological carrier and lawful category bridge are fixed.
\end{theorem}
\begin{proof}
Moise's theorem supplies standard category discipline in dimension 3: topological 3-manifolds admit compatible PL/smooth structures with the uniqueness properties required for category translation \cite{Moise1977}. This is not used as a Poincare proof; it prevents category drift from masquerading as a negative endpoint branch. Thus category variation supplies no independent negative route.
\end{proof}

\subsection{L8: Geometric or analytic auxiliary carrier}
\begin{theorem}[No unlicensed geometric carrier]
Metric, curvature, Ricci-flow, hyperbolic, spherical, analytic, or geometric-decomposition data cannot resolve the Poincare endpoint negatively unless bridged back to the fixed topological carrier.
\end{theorem}
\begin{proof}
Geometric data live in an auxiliary carrier. They may become proof content if a bridge establishes their topological significance for the same \(M\). Without such a bridge, reporting geometric status as spherehood or non-spherehood is carrier transfer. With such a bridge, the bridge content is load-bearing and routes into the relevant locus or positive bridge completion.
\end{proof}

\subsection{L9: Modification and repair}
\begin{theorem}[No repair-based negative route]
No surgery, repair, collapse, flow stage, or modification route can maintain \(\PCNeg(M)\) as a same-domain endpoint report unless carrier preservation is certified in advance.
\end{theorem}
\begin{proof}
If the operation produces a different object, the report is not about \(M\). If the operation preserves \(M\), the preservation proof and invariant it preserves are load-bearing and enter another locus. If it validates a failed or missing endpoint report after the fact, it violates act-time finality and no post-selection repair.
\end{proof}

\subsection{L10: Representative selection}
\begin{theorem}[No selector-based negative route]
A preferred triangulation, handlebody, metric, embedding, finite presentation, normal form, Heegaard splitting, or representative cannot decide negative endpoint status on the fixed carrier.
\end{theorem}
\begin{proof}
A selector chooses a representative or presentation. The endpoint readout is invariant under lawful topological redescription. Representative preference cannot perform endpoint boundary work. If a selector is accompanied by proof that its selected representative encodes an invariant endpoint route, that proof content, not the selector, is load-bearing and enters another locus. Hence selector-based negative resolution collapses.
\end{proof}

\subsection{L11: Endpoint classifier and negative-status governance}
\begin{theorem}[L11 collapse]
Endpoint-classifier or negative-status governance collapses by the no-independent-discriminator theorem.
\end{theorem}
\begin{proof}
If the classifier agrees extensionally with the declared bridge, it is bookkeeping or positive bridge completion. If it disagrees while governing the same endpoint, it is independent same-domain endpoint authority. The latter is exactly \(\IndependentSphereDisc(M)\) or a rival classifier form thereof, excluded by Theorem \ref{thm:no-independent-sphere-disc}.
\end{proof}

\begin{theorem}[No surviving endpoint-resolving negative route]\label{thm:no-surviving-negative-route}
No same-carrier endpoint-resolving negative route for the Poincare endpoint survives the route-locus audit.
\end{theorem}
\begin{proof}
By Theorem \ref{thm:route-locus-exhaustion}, every such route enters L1--L11. L1--L10 collapse by the preceding locus theorems. L11 collapses by no-independent-discriminator closure. If a route enters no locus, it does no endpoint work. Therefore no endpoint-resolving negative route survives.
\end{proof}

\section{Traditional Ricci-Flow Bridge Audit}\label{sec:ricci-bridge-audit}

\begin{mdframed}[backgroundcolor=gray!8,linewidth=0.6pt]
\textbf{Comparison boundary.} This section is not part of the proof spine. The AASC proof has already closed the endpoint by negative-branch exclusion. Ricci flow with surgery is audited only as a traditional bridge-construction route.
\end{mdframed}

The traditional route begins with a closed 3-manifold carrier, introduces a Riemannian metric as auxiliary geometric structure, evolves that structure by Ricci flow, controls singularities, performs surgery under continuation constraints, proves non-collapsing and canonical-neighborhood control, obtains finite extinction or geometrization consequences, and translates back to the topological endpoint \cite{Hamilton1982,Perelman2002,Perelman2003Surgery,Perelman2003Extinction,MorganTian2007,KleinerLott2008}.

\begin{longtable}{L{0.34\linewidth}L{0.56\linewidth}}
\toprule
\textbf{Ricci-flow component} & \textbf{AASC role}\\
\midrule
Closed 3-manifold \(M\) & Fixed topological carrier.\\
Riemannian metric & Auxiliary geometric carrier; not a replacement for \(M\).\\
Ricci flow & Bridge dynamics from geometric carrier toward classification pressure.\\
Singularity formation & Obstruction exposure inside the bridge route.\\
Surgery & Controlled continuation if carrier-preserving; not post-selection repair.\\
Non-collapsing & Bridge-admissibility control.\\
Canonical neighborhoods & Local structure control within the bridge route.\\
Finite extinction / geometrization pressure & Residue-pressure mechanism inside the traditional bridge.\\
Topological recovery & Return from auxiliary geometry to \(M\Homeo\Sthree\).\\
\bottomrule
\end{longtable}

\begin{theorem}[Ricci-flow route as bridge construction]
When its analytic and geometric burdens are discharged, the Hamilton--Perelman route is classified under UEAP as a target-preserving bridge-construction proof.
\end{theorem}
\begin{proof}
The route starts at the fixed topological carrier, introduces auxiliary metric data, controls evolution and surgery, and returns to topological classification. UEAP licenses this as bridge construction only when the auxiliary carrier is not silently substituted for \(M\), surgery is controlled as continuation rather than repair, and the final report is the original endpoint. This manuscript does not use that theorem as a premise; it classifies the traditional route after the AASC endpoint is complete.
\end{proof}

\section{Proof-Efficiency Comparison}\label{sec:efficiency}

The traditional route carries a geometric bridge burden: metric choice, flow existence, singularity control, surgery, non-collapsing, canonical neighborhoods, finite extinction or geometrization consequences, and topological recovery. The AASC route carries an endpoint-governance burden: endpoint adequacy, kernel instantiation, negative normal form, bridge-exclusion correspondence, use-kind classification, no hidden fifth case, independent-discriminator exclusion, and native negative branch exclusion.

\begin{longtable}{L{0.26\linewidth}L{0.32\linewidth}L{0.32\linewidth}}
\toprule
\textbf{Efficiency measure} & \textbf{Ricci-flow route} & \textbf{AASC route}\\
\midrule
Primary proof object & Geometric evolution of \(M\) through auxiliary metric carrier. & Fixed endpoint carrier \(\CarrierPC(M)\).\\
Main burden & Control flow, singularities, surgery, and topological recovery. & Exclude endpoint-resolving negative branch.\\
Auxiliary carrier & Riemannian metric, flow, surgery data. & Bridge slot and endpoint-status audit.\\
Failure modes & Uncontrolled singularities, surgery drift, collapse failure, recovery failure. & Kernel failure, bridge-correspondence failure, hidden fifth case, independent-discriminator survival.\\
Endpoint transition & Topology \(\to\) geometry \(\to\) topology. & Standing \(\to\) bridge exclusion audit \(\to\) native negative branch exclusion.\\
Adoption status now & Institutionally familiar. & New proof formalism.\\
\bottomrule
\end{longtable}

\begin{theorem}[AASC proof-efficiency theorem]
Relative to the Poincare endpoint, the AASC endpoint-governance route replaces geometric-evolution bridge construction with finite endpoint-use classification and independent-discriminator exclusion.
\end{theorem}
\begin{proof}
The traditional route constructs and controls a geometric evolution. The AASC route does not. It fixes the endpoint carrier, shows endpoint use forces the kernel, identifies the native negative branch with sphere-bridge exclusion, classifies official negative use as endpoint governance, routes that governance to independent sphere discrimination, and excludes such discrimination by the kernel. The endpoint transition is therefore finite at the proof-spine level.
\end{proof}

\section{Hostile-Referee Audit}\label{sec:hostile}

\subsection{Wrong-mode objections}
\begin{longtable}{L{0.44\linewidth}L{0.46\linewidth}}
\toprule
\textbf{Objection} & \textbf{Status}\\
\midrule
AASC is not mathematics. & Wrong-mode unless the critic identifies a failed definition, theorem, inference, kernel dependency, or formal constraint. AASC is the mathematical constraint formalism of this proof.\\
Lean is not a full proof of this Poincare manuscript. & Granted but irrelevant if used against the manuscript proof. Lean is audit support; the manuscript is the proof. The current Lean audit layer is available, and the stable paper-specific formalization scope includes the relevant AASC machinery directly.\\
This does not use Ricci flow. & Wrong-mode. Ricci flow is a different bridge-construction proof class.\\
This does not reproduce Perelman. & Wrong-mode. Perelman is a comparison source, not a premise.\\
This is not conventional topology. & Correct but not refuting; the proof class is AASC mathematical fixed-carrier endpoint closure.\\
A negative expression could be ordinary mathematical data. & Granted as proof support or bookkeeping. If endpoint-resolving, it enters the use-kind classification.\\
A hidden invariant might distinguish the manifold. & Then, if endpoint-resolving, it enters route-locus exhaustion or endpoint governance. If it does no endpoint work, it is not a counterexample branch.\\
A weaker kernel packet might allow negative governance. & Right-mode, not wrong-mode. The critic must define the weakening and meet the weakening-resistance burden in Section \ref{sec:weakening-resistance}.\\
\bottomrule
\end{longtable}

\subsection{Right-mode objections}
A right-mode objection must target a named theorem link:
\begin{enumerate}[label=\arabic*.]
\item AASC mathematical-status claim, by identifying a precise failure of a definition, theorem, inference, kernel dependency, or formal constraint.
\item Endpoint-under-audit declaration as target fixation rather than assumed spherehood.
\item Pointwise endpoint-use binding for arbitrary fixed \(M\).
\item Kernel-neutral endpoint adequacy.
\item Endpoint adequacy forcing \(\KPC\).
\item Poincare carrier instantiation and carrier adequacy.
\item Weakening-resistance of K5, K6, K11, and K13; specifically, whether a strict same-carrier weakening preserves target-slot fixation and minimal report evaluability while permitting independent negative governance.
\item Native negative normal form \(\PCNeg\).
\item Standard negative branch as sphere-bridge image exclusion.
\item Endpoint-used bridge exclusion as theorem-level sphere-status discrimination.
\item ATS/UEAP classification of negative branch use.
\item No endpoint-resolving non-governance.
\item Non-explosiveness of the negative-branch exclusion.
\item \(\PCNeg\) as pointwise negative endpoint branch.
\item Official negative resolution as endpoint-status governance.
\item Endpoint governance as independent sphere discrimination.
\item No independent same-domain Poincare sphere discriminator.
\item Exclusion of \(\PCNeg(M)\).
\item Sphere-bridge object \(\BridgeObjPC(M)\), including non-circularity and bridge-complete readout.
\item Local exact-complement annotation and discharge for \([\neg\SphereRead(M)]_i\).
\item Final official endpoint correspondence.
\item Lean-to-manuscript correspondence only where Lean is explicitly cited as audit support. A challenge here must identify a mismatch between the Lean-support artifact and the manuscript theorem chain; it is not enough to argue that Lean is not itself the proof.
\end{enumerate}

\subsection{Referee burden worksheet}
\begin{longtable}{L{0.32\linewidth}L{0.58\linewidth}}
\toprule
\textbf{Target} & \textbf{Successful objection must produce}\\
\midrule
AASC mathematical status & A precise failure of a definition, theorem, inference, kernel dependency, or formal constraint; not merely the assertion that AASC is philosophical, unfamiliar, or nonstandard.\\
Lean support boundary & A specific mismatch between any Lean-audited routing claim and the manuscript's stated theorem chain, if Lean is being invoked as support.\\
Manuscript proof status & A failed theorem link in the manuscript proof chain.\\
Endpoint under audit & Show this assumes \(M\Homeo\Sthree\) rather than fixing the official theorem target.\\
Pointwise endpoint use & Show an arbitrary \(M\) under the universal theorem is not under pointwise endpoint audit.\\
Carrier instantiation & Produce a closed connected simply connected 3-manifold not represented by \(\CarrierPC(M)\).\\
Native negative normal form & Produce a Poincare counterexample not expressible as \(\StandPC(M)\wedge\neg\SphereRead(M)\).\\
Bridge correspondence & Show native counterexample normal form is not sphere-bridge image exclusion.\\
Theorem-level discriminator & Show official endpoint-bearing bridge exclusion does not classify endpoint status.\\
Use-kind classification & Produce endpoint-resolving non-governance that is neither tensor, proof support, bookkeeping, nor carrier shift.\\
Non-explosiveness & Show the proof globally forbids negative mathematics rather than excluding only independent endpoint governance.\\
Weakening-resistance & Define a weaker regime \(W\) preserving fixed carrier, theorem-bearing target-slot fixation, minimal report evaluability, standing--admissibility discipline, and no carrier shift while permitting endpoint-resolving negative governance that is not proof support, bridge completion, lawful coequal target role, bookkeeping, or second endpoint authority.\\
No fifth case & Inhabit a fifth same-carrier endpoint-resolving negative occupation.\\
Independent discriminator & Show endpoint-resolving negative governance is bridge completion, proof support, carrier shift, bookkeeping, or not second endpoint authority.\\
No independent discriminator & Show independent same-domain sphere-status governance is not excluded by the kernel.\\
Strict kernel weakening & Show that a strict weakening of K5, K6, K11, or K13 preserves non-degenerate same-carrier endpoint use while allowing the negative endpoint governance excluded by Theorem \ref{thm:no-strict-weakening}.\\
Sphere-bridge object & Show that a clause of \(\BridgeObjPC(M)\) smuggles \(\SphereRead(M)\), changes carrier, or fails to preserve the official Poincare endpoint interface.\\
Bridge-complete readout & Produce a completed sphere-bridge object that does not imply \(M\Homeo\Sthree\).\\
Local reductio annotation & Prove that the exact-complement annotation is an added object-level premise rather than proof-act metadata.\\
Local countercase projection & Produce fixed-standing non-spherehood that blocks sphere-readout while not projecting to local sphere-bridge image exclusion.\\
Annotation discharge & Show that the local contradiction discharges only a strengthened premise rather than \([\neg\SphereRead(M)]_i\).\\
Official correspondence & Show \(\neg\PCNeg(M)\) does not imply \(M\Homeo\Sthree\) under the fixed native normal form.\\
\bottomrule
\end{longtable}

\section{Conclusion}\label{sec:conclusion}

The proof is a mathematical constraint-formalism proof in AASC. It is not a Ricci-flow proof, a surgery proof, or a conventional 3-manifold construction. The manuscript carries the proof: endpoint use fixes the carrier, endpoint adequacy forces the kernel, the native negative branch routes through sphere-bridge image exclusion, official negative endpoint use induces endpoint-status governance, endpoint-status governance induces independent same-domain sphere discrimination, and the kernel excludes such discrimination. The proof does not infer an admissible discriminator from true non-spherehood, and it does not forbid negative mathematics globally. It excludes the narrower forbidden structure of same-carrier endpoint-resolving negative governance outside the lawful sphere bridge. The weakening-resistance audit distinguishes target-slot fixation and minimal report evaluability from the stronger K11 and K13 closures, then shows that a weaker local packet cannot preserve non-degenerate endpoint use while permitting that forbidden structure: the alleged weakening is endpoint-equivalent, support-level, carrier-shifting, no-overreporting failure, endpoint-exhaustion failure, or second-governing. Therefore the native negative branch is impossible, and the ordinary Poincare endpoint follows.

Lean 4 support, where referenced, audits formal routing and AASC kernel discipline; it does not replace the manuscript proof. The current Poincare Lean audit archive records the endpoint-spine support surface, and the stable paper-specific formalization scope includes the relevant AASC machinery directly, while existing corpus formalization support records the maturity of the machinery used here. A same-mode refutation must therefore break the constraint chain or the Poincare instantiation, not dismiss AASC as non-mathematical language and not demand a different proof class.

The Ricci-flow comparison does not subordinate this proof to the traditional solution. It classifies the Hamilton--Perelman solution as a legitimate geometric-evolution bridge route under UEAP after the AASC endpoint has been reached. The AASC route is different: fixed endpoint use, bridge-exclusion audit, no hidden fifth case, no independent discriminator, and native negative branch exclusion. The final closeout now uses local exact-complement discharge: the assumption \([\neg\SphereRead(M)]_i\) is routed through local bridge-image exclusion and contradicted, so the discharged result is sphere-readout itself rather than a strengthened endpoint-governance premise.

\appendix

\section{Theorem Ladder}
\begin{longtable}{L{0.08\linewidth}L{0.42\linewidth}L{0.40\linewidth}}
\toprule
\textbf{No.} & \textbf{Theorem node} & \textbf{Function}\\
\midrule
0 & AASC mathematical-status lock & Fixes AASC as mathematical constraint formalism and Lean as audit support, not proof replacement.\\
1 & \(\PCEndpoint\) & Fixes official universal target, not conclusion.\\
2 & Pointwise endpoint binding & Places arbitrary \(M\) under official endpoint use and carrier instantiation.\\
3 & Kernel-neutral adequacy & Supplies target determinacy, step evaluability, act-time finality, and same-regime fidelity.\\
4 & Endpoint adequacy forces \(\KPC\) & Derives kernel roles from endpoint use.\\
5 & K1--K13 & Local fixed-carrier constraints.\\
5a & Weakening-resistance & Shows K5, K6, K11, and K13 are locally necessary against strict same-carrier weakenings that would permit independent negative governance.\\
5b & Plural routes versus plural interiors & Permits ordinary mathematical plurality while excluding plural endpoint-governing admissible interiors for one fixed endpoint slot.\\
6 & Carrier adequacy & Verifies standard Poincare data instantiate \(\CarrierPC(M)\).\\
7 & \(\PCNeg\) & Native counterexample normal form.\\
8 & \(\PCNeg\leftrightarrow\StdPCNeg\) & Removes AASC vocabulary from native branch.\\
9 & \(\StdPCNeg\leftrightarrow\SphereBridgeImgExcl\) & Identifies standard negative branch with bridge-image exclusion.\\
10 & Endpoint-used bridge exclusion & Produces theorem-level sphere-status discrimination.\\
11 & ATS/UEAP classification & Separates support, tensor endpoint work, carrier shift, bookkeeping, hidden fifth.\\
12 & No hidden fifth & Blocks endpoint-resolving non-governance.\\
12a & Non-explosiveness & Confirms negative results are not globally forbidden and isolates the forbidden Poincare structure.\\
13 & Official negative resolution & Makes endpoint-bearing \(\PCNeg\) official negative branch.\\
14 & Endpoint governance & Classifies official negative resolution as endpoint-status governance.\\
15 & Independent discriminator & Routes endpoint governance to independent sphere discrimination.\\
16 & No independent discriminator & Excludes independent same-domain sphere-status governance.\\
17 & Negative branch exclusion & Proves \(\neg\PCNeg(M)\).\\
18 & Poincare closure & Concludes \(M\Homeo\Sthree\).\\
18a & Sphere-bridge object & Fixes same-manifold, carrier, standing, readout-slot, category, UEAP, and ATS bridge discipline without containing spherehood.\\
18b & Bridge-complete sphere object & Makes bridge-completion output \(\SphereRead(M)\) explicit.\\
18c & Local reductio annotation & Records \([\neg\SphereRead(M)]_i\) as proof-role metadata, not an added premise.\\
18d & Local countercase route & Sends local non-spherehood through endpoint counterforce, local bridge-image exclusion, independent discriminator, and contradiction.\\
18e & Annotation discharge & Discharges the local assumption itself and yields sphere-readout.\\
19 & Route-locus exhaustion & Classifies endpoint-resolving negative routes.\\
20 & Ricci bridge audit & Places traditional solution in comparison layer only.\\
\bottomrule
\end{longtable}

\section{Anti-Circularity Audit}
\begin{longtable}{L{0.35\linewidth}L{0.55\linewidth}}
\toprule
\textbf{Potential circle} & \textbf{Audit response}\\
\midrule
AASC is merely philosophical & Blocked by mathematical-status lock: AASC is the mathematical constraint formalism of this proof.\\
Lean is substituted for the proof & No. The manuscript is the proof; Lean is audit support for formal routing. The current Lean audit layer is available, and the stable paper-specific formalization scope includes the relevant AASC machinery directly.\\
Endpoint under audit assumes Poincare & It fixes the official theorem target; it does not assert \(\SphereRead(M)\).\\
\(\PCNeg\) is retyped by fiat & No. It first appears as native normal form, then bridge-image exclusion, then endpoint governance only under official negative use.\\
Bridge slot contains spherehood & No. The slot fixes transition type only; completion is the endpoint result.\\
Negative branch is banned & No. Proof-support negative data are allowed. Endpoint-resolving negative use is routed and excluded after it induces independent endpoint governance.\\
Non-homeomorphism becomes its own discriminator & No. The proof does not infer a discriminator from non-homeomorphism. It infers endpoint governance from official negative endpoint use.\\
No-supported-residue implies spherehood & Not used as the final proof hinge. The final proof uses \(\neg\PCNeg(M)\Rightarrow M\Homeo\Sthree\).\\
Coordinate exhaustion proves topology by taxonomy & No. Route-locus exhaustion classifies endpoint-resolving negative use, not every possible topological fact.\\
Moise/category discipline proves Poincare & No. It prevents category drift only.\\
Ricci flow is used implicitly & No. Ricci flow appears only in comparison and bridge-audit sections.\\
The proof bans all negative results & No. The non-explosiveness theorem separates lawful negative results, proof support, carrier diagnosis, and forbidden independent endpoint governance.\\
The full K-packet is stronger than core adequacy & The weakening-resistance section audits this node by node. It distinguishes target-slot fixation and minimal report evaluability from the stronger K11 and K13 closures, then proves K5, K6, K11, and K13 as local necessities for avoiding endpoint degeneracy under same-carrier theorem-bearing use.\\
A weaker formalism could allow negative governance & A weaker formalism may allow proof support, lawful coequal negative target roles, or carrier-shift diagnosis. It cannot allow same-carrier endpoint-resolving negative governance without producing second endpoint authority or third-status degeneracy.\\
K5 bans ordinary mathematical plurality & No. K5 bans only plural endpoint-governing admissible interiors for the same fixed endpoint slot. It does not ban multiple proof routes, invariants, presentations, metrics, or support structures.\\
K13 bans uncertainty & No. It classifies uncertainty as epistemic or report status, not endpoint truth occupation.\\
\(\BridgeObjPC(M)\) contains spherehood & No. The bridge object fixes same-manifold, carrier, standing, readout-slot, category, UEAP, and ATS discipline. Spherehood appears only under bridge completion or negative-branch exclusion.\\
Local reductio assumes official negative resolution & No. The local assumption \([\neg\SphereRead(M)]_i\) is opened as exact-complement countercase. It receives local endpoint counterforce, not global official negative endpoint status.\\
Contradiction discharges a strengthened premise & No. The annotation is proof-role metadata; the only discharged object-level assumption is \(\neg\SphereRead(M)\).\\
Non-spherehood becomes a discriminator by fiat & No. It becomes local bridge-image exclusion only after fixed standing, carrier, endpoint counterforce, and bridge-role projection are in force.\\
Target-slot determinacy smuggles in K13 & No. Target-slot determinacy fixes what is evaluated. K13 is the separate fail-closed theorem showing that unknown, deferred, unsupported, bookkeeping, and hidden-fifth statuses cannot occupy endpoint truth.\\
Minimal report evaluability smuggles in K11 & No. Minimal report evaluability only attaches reports to the fixed target and classifies them by role. K11 is the separate no-overreporting theorem that endpoint force may not exceed carrier, bridge, and tensor support.\\
\bottomrule
\end{longtable}

\section{Preferred Matrix Source-Use Audit}
\begin{longtable}{L{0.30\linewidth}L{0.30\linewidth}L{0.30\linewidth}}
\toprule
\textbf{Local theorem} & \textbf{Preferred corpus support} & \textbf{Role}\\
\midrule
AASC mathematical status & AASC kernel corpus and formal support surface & Treats AASC as mathematical constraint formalism, not philosophical commentary.\\
Lean support boundary & Project formalization discipline & Keeps Lean as audit support and identifies the paper-specific formalization scope, not replacement proof.\\
Endpoint adequacy & Kernel of Admissibility & Shows endpoint use requires reference, standing, admissibility, and irreversibility.\\
Weakening resistance & Kernel of Admissibility; Standing-Admissibility Identity; UEAP; Bivalence & Shows strict weakenings of K5, K6, K11, and K13 either preserve the endpoint packet extensionally or exit non-degenerate same-carrier endpoint use; distinguishes minimal evaluability from K11 no-overreporting and target-slot fixation from K13 exhaustion.\\
Bivalence / fail-closed status & Bivalence / Structure of Admissibility & Blocks third endpoint truth statuses.\\
Unique admissible interior & Structure of Admissibility & Blocks plural endpoint-status authorities.\\
Standing--admissibility identity & Standing-Admissibility Identity Closure & Blocks standing outside the admissibility gate.\\
No generators & Impossibility Suite / Bivalence & Endpoint labels cannot generate standing.\\
No carriers & Bivalence / AMetric Boundary & Surrogate data cannot carry endpoint truth without bridge.\\
No repairs & AMetric Boundary / Kernel & Later data cannot repair failed same-act reports.\\
AMetric boundary & AMetric Boundary & Blocks selector authority.\\
UEAP report preservation & Claim Standing and Legitimacy & Reports cannot exceed support.\\
ATS layer discipline & Anchor, Tensor, and Skin & Skin, labels, and bookkeeping cannot become tensor proof.\\
No hidden fifth case & ATS / UEAP / endpoint exhaustion & Endpoint resolution without governance is impossible.\\
No independent discriminator & Standing-admissibility identity / unique interior & Blocks second same-domain endpoint authority.\\
\bottomrule
\end{longtable}

\section{Faithful Same-Regime Counterexample Worksheet}
A proposed same-mode counterexample or refutation must complete the following worksheet. Failure to complete it means the proposal is not a same-mode refutation of the proof spine.

\begin{longtable}{L{0.34\linewidth}L{0.56\linewidth}}
\toprule
\textbf{Target} & \textbf{Successful objection must produce}\\
\midrule
AASC mathematical status & Identify a precise failure of a definition, theorem, inference, kernel dependency, or formal constraint.\\
Lean support boundary & Identify a mismatch between any Lean-supported audit claim and the manuscript theorem chain, if Lean is invoked.\\
Manuscript proof status & Identify a failed theorem link in the proof chain stated in this manuscript.\\
Endpoint-under-audit & Show target fixation assumes the conclusion.\\
Pointwise endpoint use & Show the arbitrary fixed \(M\) is not under pointwise endpoint audit.\\
Carrier adequacy & Identify ordinary Poincare data not represented by \(\CarrierPC(M)\).\\
Kernel denial & Preserve theorem-bearing endpoint use while abandoning a named kernel role.\\
Native normal form & Produce a counterexample not expressible as \(\StandPC(M)\wedge M\not\Homeo\Sthree\).\\
Bridge correspondence & Show \(\PCNeg\) is not sphere-bridge image exclusion.\\
Endpoint-status classification & Show official bridge exclusion does not classify the fixed endpoint.\\
No hidden fifth & Exhibit endpoint-resolving non-governance.\\
Independent discriminator & Show negative endpoint governance that is not second endpoint authority.\\
No-independent-discriminator theorem & Show such second authority is compatible with \(\KPC\).\\
Final correspondence & Show \(\neg\PCNeg(M)\) fails to yield \(M\Homeo\Sthree\).\\
\bottomrule
\end{longtable}


\section{Lean/AASC Support Appendix and Formalization Scope}

\begin{mdframed}[backgroundcolor=gray!8,linewidth=0.6pt]
\textbf{Support boundary.} The proof of this manuscript is the mathematical theorem chain above. The Lean 4 layer is support/audit material for the proof-spine routing and AASC kernel discipline; it is not substituted for the proof text. The public Poincare Lean audit archive records the manuscript-facing endpoint route, the reusable AASC foundation layer, focused axiom checks, the pre-Lean manuscript signature map, and the manuscript-facing PDF/source snapshot \cite{PoincareLeanAuditGithub2026}. The Zenodo DOI record supplies a citable archival pointer for the same Lean audit layer \cite{PoincareLeanAuditZenodo2026}. The stable paper-specific formalization scope includes the relevant AASC machinery directly: kernel discipline, below-kernel underivability posture, endpoint routing, no-hidden-fifth closure, and no-independent-classifier structure.

The explicit sphere-bridge object and local exact-complement reductio annotation introduced in this release are manuscript-level defensive refinements of the already represented endpoint route. They do not add a new first-principles topology or Ricci-flow formalization burden; they clarify the same carrier-preserving bridge role and the local assumption-discharge discipline used by the manuscript closeout.
\end{mdframed}

\subsection*{Public audit record}
\phantomsection
\addcontentsline{toc}{subsection}{Public audit record}

The current public Lean support layer is the standalone repository \emph{AASC Poincare Endpoint Lean Audit}, release \texttt{v1.0.1}. The repository describes itself as a Lean 4 archive for the AASC-first Poincare endpoint proof spine and records a complete AASC endpoint-structure proof route for the Poincare endpoint in the audit-spine proof class \cite{PoincareLeanAuditGithub2026}. Its README states the manuscript-facing route: official fixed-carrier Poincare endpoint use enters the kernel-governed endpoint regime; \texttt{PCNeg} is identified with standard negative form and sphere-bridge image exclusion; official negative endpoint use induces theorem-level same-domain endpoint-status governance; that governance induces an independent sphere discriminator; and the reusable AASC no-independent-classifier closure excludes the discriminator, forcing pointwise \texttt{SphereRead} and yielding \texttt{OfficialPCEndpoint} \cite{PoincareLeanAuditGithub2026}.

\begin{longtable}{L{0.30\linewidth}L{0.60\linewidth}}
\toprule
\textbf{Audit item} & \textbf{Recorded support}\\
\midrule
Public repository & \emph{AASC-Poincare-Endpoint-Lean-Audit}; see \cite{PoincareLeanAuditGithub2026}.\\
Current release & \texttt{v1.0.1}, \emph{AASC Poincare Endpoint Lean Audit v1.0.1}.\\
Archival DOI & \texttt{10.5281/zenodo.20620926}; see \cite{PoincareLeanAuditZenodo2026}.\\
Repository role & Standalone Lean 4 audit archive for the AASC-first Poincare endpoint proof spine; separated from the broader AASC working checkout.\\
Included layers & Reusable AASC foundation layer; Poincare endpoint Lean audit module; focused axiom checks; pre-Lean manuscript signature map; manuscript-facing PDF/source snapshot.\\
Proof boundary & Audit support for the endpoint-structure route; not a replacement for the manuscript proof and not a conventional first-principles topology or Ricci-flow formalization.\\
\bottomrule
\end{longtable}

\subsection*{Verification commands and audit status}
\phantomsection
\addcontentsline{toc}{subsection}{Verification commands and audit status}

The repository records the following verification commands for the focused Poincare endpoint audit surface \cite{PoincareLeanAuditGithub2026,PoincareLeanAuditRelease2026}:
\begin{verbatim}
lake build MaleyLean.Papers.Poincare.AuditRunners
powershell -ExecutionPolicy Bypass -File \
  scripts/check-poincare-endpoint-audit.ps1
\end{verbatim}

\begin{longtable}{L{0.38\linewidth}L{0.52\linewidth}}
\toprule
\textbf{Audit status entry} & \textbf{Recorded status}\\
\midrule
Endpoint closure & \texttt{PCEndpointClosure=100\%}.\\
Ricci-flow comparison boundary & \texttt{PCRicciBridgeComparisonBoundary=100\%}.\\
Reference archive maturity & 100\% comparable maturity for the Poincare reference archive.\\
GitHub Actions & Audit passed for release \texttt{v1.0.1}.\\
Active audit surface & No live project-level \texttt{axiom}, \texttt{sorry}, \texttt{admit}, or \texttt{unsafe} declaration in the active Poincare audit surface.\\
Pre-Lean signature map & Parses outside the active proof surface and is not imported by the active Lean audit surface.\\
Ricci-flow material & Comparison boundary only; not proof machinery.\\
\bottomrule
\end{longtable}

\subsection*{Pinned environment and release assets}
\phantomsection
\addcontentsline{toc}{subsection}{Pinned environment and release assets}

The GitHub release records the pinned Lean environment and release assets as follows \cite{PoincareLeanAuditRelease2026}.

\begin{longtable}{L{0.32\linewidth}L{0.58\linewidth}}
\toprule
\textbf{Field} & \textbf{Value}\\
\midrule
Lean toolchain & \texttt{leanprover/lean4:v4.28.0}.\\
mathlib revision & \begin{tabular}[t]{@{}l@{}}\texttt{8f9d9cff6bd728b17a24e163}\
\texttt{c9402775d9e6a365}\end{tabular}.\\
Release assets & Standalone archive zip, manifest, and SHA256 checksum file.\\
Release archive & The release handoff records that the attached \texttt{.sha256} file stores the archive checksum.\\
\bottomrule
\end{longtable}

\subsection*{Lean handoff chain}
\phantomsection
\addcontentsline{toc}{subsection}{Lean handoff chain}

The audit handoff records the closed endpoint route in Lean-facing form \cite{PoincareLeanAuditHandoff2026}:
\begin{verbatim}
OfficialPCEndpointUse(M)
  -> PCKernelInstantiated(M)
  -> PCNeg(M) as SphereBridgeImgExclPC(M)
  -> ThmSphereDiscPC(M)
  -> EndpointGovPC(M)
  -> IndependentSphereDiscPC(M)
  -> contradiction
  -> SphereRead(M)
  -> OfficialPCEndpoint
\end{verbatim}
The local closeout refinement has the parallel proof-role shape
\begin{verbatim}
[not SphereRead(M)]_i
  -> LocalCountercasePC(M)
  -> EndpointCounterforcePC(M)
  -> LocalSphereBridgeImageExclusion(M)
  -> IndependentSphereDiscPC(M)
  -> contradiction
  -> SphereRead(M)
\end{verbatim}
This is the Lean-facing counterpart of the manuscript proof spine. The handoff also records that Ricci-flow material is represented only by named comparison-boundary objects and is not active proof machinery \cite{PoincareLeanAuditHandoff2026}.

\subsection*{Principal Lean anchors}
\phantomsection
\addcontentsline{toc}{subsection}{Principal Lean anchors}

The README and handoff identify the following principal anchors for the Poincare endpoint audit surface \cite{PoincareLeanAuditGithub2026,PoincareLeanAuditHandoff2026}. The names are shown in grouped form to preserve line breaking while retaining exact audit traceability.

\begin{longtable}{L{0.32\linewidth}L{0.58\linewidth}}
\toprule
\textbf{Anchor group} & \textbf{Recorded Lean anchors and manuscript role}\\
\midrule
Endpoint module & \texttt{EndpointClosure.lean}: main Poincare endpoint closure module under \texttt{MaleyLean/Papers/Poincare/}.\\
Negative occupation and bridge exclusion & \begin{tabular}[t]{@{}l@{}}\texttt{pcNegativeOccupation\_nonoptional}\\\texttt{pcNeg\_iff\_sphereBridgeImageExclusion}\\\texttt{pcNeg\_impossible}\end{tabular}\\[-0.2em]
& These anchors record nonoptional negative occupation, the bridge-exclusion correspondence, and exclusion of the native negative branch.\\
Endpoint governance and hidden fifth case & \begin{tabular}[t]{@{}l@{}}\texttt{pcOfficialNegativeResolution\_}\\\texttt{endpointStatusGovernance}\\\texttt{pcEndpointResolvingNonGovernance\_}\\\texttt{hiddenFifthCase\_impossible}\end{tabular}\\[-0.2em]
& These anchors record governance from official negative resolution and impossibility of endpoint-resolving non-governance.\\
Independent discriminator closure & \begin{tabular}[t]{@{}l@{}}\texttt{pcNeg\_independentSphereDiscriminator}\\\texttt{pcNoIndependentSphereDiscriminator\_}\\\texttt{of\_foundationalNoClassifier}\end{tabular}\\[-0.2em]
& These anchors route negative endpoint governance to independent sphere discrimination and close it by no-classifier support.\\
Endpoint closeout & \begin{tabular}[t]{@{}l@{}}\texttt{sphereRead\_forced}\\\texttt{officialPCEndpoint\_of\_aascContext}\\\texttt{pcEndpointAASCContext\_closes\_endpoint}\end{tabular}\\[-0.2em]
& These anchors record forced pointwise sphere-readout and the official Poincare endpoint closeout.\\
Ricci-flow comparison boundary & \begin{tabular}[t]{@{}l@{}}\texttt{ricciBridgeAudit\_context\_iff}\\\texttt{pcRicciBridge\_remains\_comparisonOnly}\end{tabular}\\[-0.2em]
& These anchors record the Ricci-flow comparison boundary.\\
\bottomrule
\end{longtable}

\subsection*{Truth boundary}
\phantomsection
\addcontentsline{toc}{subsection}{Truth boundary}

The Lean audit archive is a support surface for the AASC endpoint-structure route. It does not claim to formalize 3-manifold topology, homeomorphism, Ricci flow, surgery, or classical topological classification from first principles; these are represented by semantic audit carriers or comparison standing so that the AASC endpoint route can be checked without importing a full external topology library \cite{PoincareLeanAuditGithub2026,PoincareLeanAuditRelease2026}. This boundary matches the proof-class lock of the manuscript: the proof class is AASC constraint-formalism endpoint closure, not a Hamilton--Perelman formalization.

\begin{longtable}{L{0.32\linewidth}L{0.58\linewidth}}
\toprule
\textbf{Support object} & \textbf{Status boundary}\\
\midrule
AASC kernel machinery & Included through the reusable AASC foundation layer and Poincare endpoint audit surface; also supported by existing corpus formalization provenance.\\
Below-kernel underivability posture & Support/audit material for the non-derivability of the kernel roles from lower same-domain resources without presupposition.\\
Poincare endpoint Lean layer & Current audit layer is public in release \texttt{v1.0.1}; the stable manuscript formalization scope includes the relevant AASC machinery directly.\\
Manuscript proof status & The manuscript theorem chain remains the proof text. The Lean layer audits the route and formal discipline.\\
Classical topology libraries & Not used as proof premises for the endpoint; classical sources remain comparison, context, or category-discipline support.\\
\bottomrule
\end{longtable}

\subsection*{Formalization map}
\phantomsection
\addcontentsline{toc}{subsection}{Formalization map}
\begin{verbatim}
PCEndpointUnderAudit
Closed3Carrier M
PoincareStanding M
SphereRead M
FixedPoincareCarrier M
OfficialPCEndpointUse M
EndpointAdeqPC M
KPC M
AASCMathematicalStatus
LeanSupportBoundary
PCSphereBridgeObject M
PCSphereBridgeObjectComplete M
PCLocalReductioCountercase M
PCLocalNegativeCountercase M
PCLocalSphereBridgeImageExclusion M
PCNeg M
StdPCNeg M
BridgeSlotPC M
BridgeAdmPC M
BridgeCompPC M
SphereBridgeImgExclPC M
OfficialPCNegativeResolution M
ThmSphereStatusDiscPC M
EndpointGovPC M
IndependentPCSphereDiscriminator M
NegativeResultsNotGloballyForbidden M
NoHiddenFifthPC M
NoIndependentSphereDiscriminator M
PCNegExcluded M
AASCPoincareClosure M
RouteLocus M L_i
RouteLocusExhaustion M
CollapseRouteLocus L_i
RicciBridgeAudit M
HostileRefereeBurden M
\end{verbatim}

\begin{thebibliography}{99}
\bibitem{AASCMatrix2026} Amos Jay Maley. \emph{AASC Corpus Control Matrix}. Internal corpus matrix, 2026.
\bibitem{PreferredMatrixScan2026} Amos Jay Maley. \emph{Poincare Preferred-Matrix Component Scan and Import Plan}. Internal scan outputs, 2026.
\bibitem{AASC_Kernel_2026} Amos Jay Maley. \emph{Non-Degenerate Construction and the Kernel of Admissibility}. Preferred hardened AASC source family, 2026.
\bibitem{AASC_Structure_Admissibility_2026} Amos Jay Maley. \emph{The Structure of Admissibility: Ametricity, Bivalence, and the Unique Admissible Interior}. Preferred hardened AASC source family, 2026.
\bibitem{AASC_StandingIdentity_2026} Amos Jay Maley. \emph{Standing-Admissibility Identity Closure: A Fixed-Domain Meta-Necessity Closure}. Preferred hardened AASC source family, 2026.
\bibitem{AASC_ClaimStanding_UEAP_2026} Amos Jay Maley. \emph{Claim Standing and Legitimacy: A Universal Audit Protocol for Laundering Exclusion and Reportable Output}. Preferred hardened AASC source family, 2026.
\bibitem{AASC_AMetric_2026} Amos Jay Maley. \emph{The AMetric Boundary as a Unique Regime-Invariant Constraint on Admissibility-Bearing Construction}. Preferred hardened AASC source family, 2026.
\bibitem{AASC_ATS_2026} Amos Jay Maley. \emph{Anchor, Tensor, and Skin: A Fixed-Domain Decomposition Theorem for Admissible Description}. Preferred hardened AASC source family, 2026.
\bibitem{AASC_Bivalence_2026} Amos Jay Maley. \emph{Bivalence Theorem for Non-Degenerate Construction}. Preferred hardened AASC source family, 2026.
\bibitem{ClayPoincare} Clay Mathematics Institute. \emph{The Poincare Conjecture}. Problem statement and solved-status context. \url{https://www.claymath.org/millennium/poincare-conjecture/}.
\bibitem{ClayMorganTian} Clay Mathematics Institute. \emph{Ricci Flow and the Poincare Conjecture, by John Morgan and Gang Tian}. \url{https://www.claymath.org/resource/ricci-flow-and-the-poincare-conjecture/}.
\bibitem{MorganTian2007} John W. Morgan and Gang Tian. \emph{Ricci Flow and the Poincare Conjecture}. Clay Mathematics Monographs, Vol. 3, AMS/CMI, 2007.
\bibitem{Hamilton1982} Richard S. Hamilton. Three-manifolds with positive Ricci curvature. \emph{Journal of Differential Geometry} 17(2):255--306, 1982.
\bibitem{Perelman2002} Grisha Perelman. \emph{The entropy formula for the Ricci flow and its geometric applications}. arXiv:math/0211159, 2002.
\bibitem{Perelman2003Surgery} Grisha Perelman. \emph{Ricci flow with surgery on three-manifolds}. arXiv:math/0303109, 2003.
\bibitem{Perelman2003Extinction} Grisha Perelman. \emph{Finite extinction time for the solutions to the Ricci flow on certain three-manifolds}. arXiv:math/0307245, 2003.
\bibitem{KleinerLott2008} Bruce Kleiner and John Lott. \emph{Notes on Perelman's Papers}. arXiv:math/0605667, 2008.
\bibitem{Moise1977} Edwin E. Moise. \emph{Geometric Topology in Dimensions 2 and 3}. Graduate Texts in Mathematics 47, Springer, 1977.

\bibitem{PoincareLeanAuditGithub2026} Amos Jay Maley. \emph{AASC Poincare Endpoint Lean Audit}. Public Lean 4 audit repository, release v1.0.1, 2026. \url{https://github.com/somamaley-ux/AASC-Poincare-Endpoint-Lean-Audit}.
\bibitem{PoincareLeanAuditRelease2026} Amos Jay Maley. \emph{AASC Poincare Endpoint Lean Audit v1.0.1}. GitHub release, 2026. \url{https://github.com/somamaley-ux/AASC-Poincare-Endpoint-Lean-Audit/releases/tag/v1.0.1}.
\bibitem{PoincareLeanAuditHandoff2026} Amos Jay Maley. \emph{Poincare Audit Handoff}. GitHub repository handoff file, 2026. \url{https://github.com/somamaley-ux/AASC-Poincare-Endpoint-Lean-Audit/blob/main/AUDIT_HANDOFF.md}.
\bibitem{PoincareLeanAuditZenodo2026} Amos Jay Maley. \emph{AASC Poincare Endpoint Lean Audit}. Zenodo archive DOI 10.5281/zenodo.20620926, 2026. \url{https://doi.org/10.5281/zenodo.20620926}.
\end{thebibliography}

\end{document}
